\chapter{Frequenzgang zeitdiskreter Systeme}\label{sec:frequenzgang-zeitdiskreter-systeme}

In Kapitel \hyperref[zeitdiskrete-systeme-im-zeitbereich]{4} und \hyperref[zeitdiskrete-systeme-im-z-bereich]{6} werden Systeme im Zeit- und Bildbereich behandelt. Es wird gezeigt, wie Systemantworten berechnet werden können. Außerdem wird das Verhalten der Systeme mit Hilfe von Übertragungsfunktionen beschrieben.

Die Fourier-Transformation für zeitkontinuierliche Signale und Systeme ermöglicht es, die Änderung des Spektrums eines Signals beim Durchlaufen eines dynamischen Systems zu beschreiben. Die Änderung kann durch die Multiplikation des Spektrums vom Eingangssignal mit dem Frequenzgang des Systems beschrieben werden. Dieser Zusammenhang wird auf verschiedene Arten auch für zeitdiskrete Systeme hergeleitet.

Der Frequenzgang eines zeitdiskreten Systems kann wie bei zeitkontinuierlichen Systemen mit Ortskurven, Frequenzgangkennlinien und Bode-Diagrammen beschrieben werden. Alle Darstellungsformen werden eingeführt, wobei sich die Darstellungen im Wesentlichen auf Frequenzgangskennlinien und Bode-Diagramme konzentrieren.

Besonders anschauliche Darstellungen von Frequenzgangkennlinien ergeben sich für Systeme mit gebrochen rationalen Übertragungsfunktionen. Für die Darstellung der zugehörigen Frequenzgangskennlinien können die Übertragungsfunktionen in Linearfaktoren zerlegt werden. Damit lässt sich der Frequenzgang eines Systems auf die Summe der logarithmierten Amplituden- und Phasengänge von Linearfaktoren zurückführen. Dieses Verfahren wird vorgestellt und an Beispielen vertieft.

Systeme mit endlicher Impulsantwort (FIR-Systeme) weisen oft eine Symmetrie auf, die das Berechnen des Frequenzgangs vereinfacht. Erfüllt die Symmetrie bestimmte Kriterien, wird der Phasengang von FIR-Systemen linear. Die Anforderungen werden analysiert und Vorteile eines linearen Phasengangs werden an Beispielen verdeutlicht.

Das Kapitel schließt mit Übungsaufgaben zum Frequenzgang zeitdiskreter Systeme. Diese Übungsaufgaben verdeutlichen den Zusammenhang zwischen Zeitbereich, Bildbereich und Frequenzbereich.

\section{Motivation und Herleitung}\label{sec:8.1}

Der Begriff des Frequenzgangs eines zeitdiskreten Systems kann auf unterschiedlichen Wegen erklärt werden. Um das Wissen aus zeitkontinuierlichen Systemen, der Systembeschreibung im Zeit- und im\\
z-Bereich sowie die Fourier-Transformation von Signalfolgen miteinander zu vernetzen, werden an dieser Stelle unterschiedliche Motivationen und Herleitungen vorgestellt.

\subsection{Berechnung des Frequenzgangs aus der Differenzengleichung eines Systems}\label{sec:8.1.1}

Ausgangspunkt für die Herleitung des Frequenzgangs eines Systems ist die Differenzengleichung
\begin{align}
\sum_{n = 0}^{N}{c_{n}  y[ k - n]} &= \sum_{m = 0}^{M}{d_{m}  u[ k - m]} \label{eq:8.1}
\end{align}
Mit der Verschiebungsregel der Fourier-Transformation für Signalfolgen kann die Gleichung in den Frequenzbereich transformiert werden:
\begin{align}
\sum_{n = 0}^{N}{c_{n}  Y(\Omega)  e^{- j  \Omega  n}} &= \sum_{m = 0}^{M}{d_{m}  U(\Omega)  e^{- j  \Omega  m}} \label{eq:8.2}
\end{align}
Durch Ausklammern von \(U(Ω)\) und \(Y(Ω)\) ergibt sich
\begin{align}
G(\Omega) &= \frac{Y(\Omega)}{U(\Omega)} = \frac{\sum_{m = 0}^{M}{d_{m}  e^{- j  \Omega  m}}}{\sum_{n = 0}^{N}{c_{n}  e^{- j  \Omega  n}}} \label{eq:8.3}
\end{align}
Die resultierende Funktion \(G(Ω)\) wird als Frequenzgang des Systems bezeichnet. In Abschnitt \ref{sec:8.1.3} wird sich zeigen, dass der Frequenzgang \(G(Ω)\) für kausale und stabile Systeme gleichzeitig die Fourier-Transformierte der Impulsantwort \(g[k]\) ist.

\begin{example}{Rekursives Tiefpass-Filter}
\label{beispiel-rekursives-tiefpass-filter}

Das Vorgehen wird am Beispiel eines rekursiven Tiefpassfilters verdeutlicht. Die Transformation der Differenzengleichung
\begin{align}
y[ k] - GF  y[ k - 1] &= (1 - GF)  u[ k] \label{eq:8.4}
\end{align}
in den Frequenzbereich führt mit 0 \textless{} GF \textless{} 1 zu der Gleichung
\begin{align}
Y(\Omega) - GF  Y(\Omega)  e^{- j  \Omega} &= (1 - GF)  U(\Omega) \label{eq:8.5}
\end{align}
Ein Auflösen der Gleichung führt zu der Übertragungsfunktion
\begin{align}
G(\Omega) &= \frac{Y(\Omega)}{U(\Omega)} = \frac{1 - GF}{1 - GF  e^{- j  \Omega}} \label{eq:8.6}
\end{align}
Der Frequenzgang wird im Abschnitt \ref{sec:8.2} dargestellt und interpretiert.

\end{example}

\subsection{Berechnung des Frequenzgangs aus der Übertragungsfunktion \(G(z)\) eines Systems}\label{sec:berechnung-des-frequenzgangs-aus-der-uxfcbertragungsfunktion-gz-eines-systems}

Die Berechnung der Übertragungsfunktion \(G(Ω)\) wird in Abschnitt \ref{sec:8.1.1} formal genauso durchgeführt wie die Berechnung der Übertragungsfunktion im z-Bereich. Ein Vergleich der Fourier- und der z-Transformierten der Differenzengleichung
\begin{align}
\sum_{n = 0}^{N}{c_{n}  Y(z)  z^{- n}} &= \sum_{m = 0}^{M}{d_{m}  U(z)  z^{- m}} \label{eq:8.7}
\end{align}
zeigt, dass die beiden Transformierten und damit auch die beiden Übertragungsfunktionen übereinstimmen, wenn die Variable z durch den Ausdruck \(e^{jΩ}\) substituiert wird. Aus Kapitel \ref{sec:7.3.2} ist bekannt, dass diese Substitution nur dann durchgeführt werden darf, wenn der Einheitskreis im Konvergenzbereich der z-Transformierten \(G(z)\) liegt. Das ist bei asymptotisch stabilen Systemen der Fall, sodass für asymptotisch stabile Systeme gilt:
\begin{align}
G(\Omega) &= \left. \ G(z) \right|_{z = e^{j  \Omega}} \label{eq:8.8}
\end{align}
\begin{example}{Rekursives Tiefpass-Filter}
\label{beispiel-rekursives-tiefpass-filter-1}

Das rekursive Tiefpass-Filter hat im z-Bereich die Übertragungsfunktion
\begin{align}
G(z) &= \frac{1 - GF}{1 - GF  z^{- 1}} = \frac{(1 - GF)  z}{z - GF} \label{eq:8.9}
\end{align}
Der Pol der Übertragungsfunktion \(G(z)\) liegt an der Stelle
\begin{align}
\alpha &= GF \label{eq:8.10}
\end{align}
Für 0 \textless{} GF \textless{} 1 liegt der Pol innerhalb des Einheitskreises. Da in diesem Fall das System stabil ist, ergibt sich für diesen Fall der Frequenzgang des Systems zu
\begin{align}
G(\Omega) &= \left. \ G(z) \right|_{z = e^{j  \Omega}} = \frac{1 - GF}{1 - GF  e^{- j  \Omega}} \label{eq:8.11}
\end{align}
\end{example}

\subsection{Faltungsregel der Fourier-Transformation von Signalfolgen}\label{sec:8.1.3}

Bei der Beschreibung von zeitdiskreten Systemen wird das Ausgangssignal \(y[k]\) eines Systems über die Faltung von Eingangssignal \(u[k]\) und Impulsantwort \(g[k]\) berechnet.
\begin{align}
y[ k] &= g[ k]*u[ k] = u[ k]*g[ k] \label{eq:8.12}
\end{align}
Mit der Faltungsregel der Fourier-Transformation ergibt sich im Frequenzbereich der Zusammenhang
\begin{align}
Y(\Omega) &= G(\Omega)  U(\Omega) = U(\Omega)  G(\Omega) \label{eq:8.13}
\end{align}
Das Spektrum \(Y(Ω)\) des Ausgangssignals \(y[k]\) errechnet sich aus dem Produkt der Spektren der Impulsantwort \(G(Ω)\) und des Spektrums des Eingangssignals \(U(Ω)\). Charakteristisch für das System ist die Übertragungsfunktion \(G(Ω)\), die sich aus dem Quotienten der Fourier-Transformierten von Aus- und Eingangssignal ergibt.
\begin{align}
G(\Omega) &= \frac{Y(\Omega)}{U(\Omega)} \label{eq:8.14}
\end{align}
Sie entspricht gemäß der obigen Herleitung der Fourier-Transformierten der Impulsantwort \(g[k]\).

\begin{example}{Rekursives Tiefpass-Filter}
\label{beispiel-rekursives-tiefpass-filter-2}

Für ein rekursives Filter mit der Differenzengleichung
\begin{align}
y[ k] - GF  y[ k - 1] &= (1 - GF)  u[ k] \label{eq:8.15}
\end{align}
ergibt sich für 0 \textless{} GF \textless{} 1 die gegen null konvergierende Impulsantwort
\begin{align}
g[ k] &= GF^{k}  \sigma[ k] \label{eq:8.16}
\end{align}
Aus der Impulsantwort wird mit den Rechenregeln der Fourier-Transformation der Frequenzgang bestimmt zu
\begin{align}
G(\Omega) &= \frac{1 - GF}{1 - GF  e^{- j  \Omega}} \label{eq:8.17}
\end{align}
Alle Wege, den Frequenzgang des rekursiven Filters zu beschreiben, führen damit zu demselben Ergebnis.

\end{example}

\subsection{Reaktion zeitdiskreter Systeme auf eine kausale, harmonische Anregung}\label{sec:reaktion-zeitdiskreter-systeme-auf-eine-kausale-harmonische-anregung}

Bei zeitkontinuierlichen, kausalen Systemen wird gezeigt, dass ein System bei Anregung mit einer kausalen harmonischen Anregung mit einem Signal antwortet, das sich aus einem Einschwingvorgang und einer harmonischen Antwort gleicher Frequenz zusammensetzt. Dieser Gedanke wird an dieser Stelle aufgegriffen und auf zeitdiskrete Systeme angewendet. Dazu wird zunächst die Systemreaktion auf ein Eingangssignal \(u_{1}\){[}k{]} der Form
\begin{align}
u_{1}[ k] &= e^{j  \Omega  k}  \sigma[ k] \label{eq:8.18}
\end{align}
berechnet. Da das System und die Eingangsfolge kausal sind, ist die Ausgangsfolge für k \textless{} 0 null. Für k ≥ 0 ergibt sich die Ausgangsfolge aus einer endlichen Faltungssumme.
\begin{align}
y_{1}[ k] &= g[ k]*u_{1}[ k] = \sum_{\kappa = - \infty}^{\infty}{g[\kappa]}  u_{1}[ k - \kappa] = \sum_{\kappa = 0}^{k}{g[\kappa]}  e^{j  \Omega  (k - \kappa)} = e^{j  \Omega  k}  \sum_{\kappa = 0}^{k}{g[\kappa]}  e^{- j  \Omega  \kappa} \label{eq:8.19}
\end{align}
Um auf die Gleichungen für die Fourier-Transformation zu kommen, wird die Summe erweitert:
\begin{align}
{y_{1} \left[k\right]} &=  {e^{j \Omega  k}  \sum _{\kappa =0}^{k}g\left[\kappa \right]  e^{-j \Omega  \kappa } =e^{j \Omega  k}  \sum _{\kappa =0}^{\infty }g\left[\kappa \right]  e^{-j \Omega  \kappa }  \sigma \left[k\right]-e^{j \Omega  k}  \sum _{\kappa =k+1}^{\infty }g\left[\kappa \right]  e^{-j \Omega  \kappa }  \sigma \left[k\right]} \\
{} &=  {e^{j \Omega  k}  G\left(\Omega \right) \sigma \left[k\right]-e^{j \Omega  k}  \sum _{\kappa =k+1}^{\infty }g\left[\kappa \right]  e^{-j \Omega  \kappa }  \sigma \left[k\right]}
 \label{eq:8.20}
\end{align}
Analog ergibt sich für die Eingangsfolge
\begin{align}
u_{2}[ k] &= e^{- j  \Omega  k}  \sigma[ k] \label{eq:8.21}
\end{align}
die Ausgangsfolge
\begin{align}
{y_{2} \left[k\right]} &=  {e^{-j \Omega  k}  \sum _{\kappa =0}^{k}g\left[n\right]  e^{j \Omega  \kappa } =e^{-j \Omega  k}  \sum _{\kappa =0}^{\infty }g\left[\kappa \right]  e^{j \Omega  \kappa }  \sigma \left[k\right]-e^{-j \Omega  k}  \sum _{\kappa =k+1}^{\infty }g\left[\kappa \right]  e^{j \Omega  \kappa }  \sigma \left[k\right]} \\
{} &=  {e^{-j \Omega  k}  G\left(-\Omega \right) \sigma \left[k\right]-e^{-j \Omega  k}  \sum _{\kappa =k+1}^{\infty }g\left[\kappa \right]  e^{j \Omega  \kappa }  \sigma \left[k\right]}
 \label{eq:8.22}
\end{align}
Mit der Eulerschen Formel kann damit die Systemreaktion auf eine kausale, harmonische Anregung der Form
\begin{align}
u[ k] &= A  \cos(\Omega  k)  \sigma[ k] = \frac{A}{2}  \left( e^{j  \Omega  k} + e^{- j  \Omega  k} \right)  \sigma[ k] \label{eq:8.23}
\end{align}
berechnet werden. Es ergibt sich die Folge
\begin{align}
{y\left[k\right]} &=  {\frac{A}{2}  \left(e^{j \Omega  k}  G\left(\Omega \right)-e^{j \Omega  k}  \sum _{\kappa =k+1}^{\infty }g\left[\kappa \right]  e^{-j \Omega  \kappa } \right) \sigma \left[k\right]} \\
{} & {} & {+\frac{A}{2}  \left(e^{-j \Omega  k}  G\left(-\Omega \right)-e^{-j \Omega  k}  \sum _{\kappa =k+1}^{\infty }g\left[\kappa \right]  e^{j \Omega  \kappa } \right) \sigma \left[k\right]} \\
{} &=  {\frac{A}{2}  \left(e^{j \Omega  k}  G\left(\Omega \right)+e^{-j \Omega  k}  G\left(-\Omega \right)\right) \sigma \left[k\right]} \\
{} & {} & {-\frac{A}{2}  \left(\sum _{n=k+1}^{\infty }g\left[n\right]  e^{j \Omega  \left(k-n\right)} +\sum _{n=k+1}^{\infty }g\left[n\right]  e^{-j \Omega  \left(k-n\right)} \right) \sigma \left[k\right]}
 \label{eq:8.24}
\end{align}
Aufgrund der Symmetrieregeln der Fourier-Transformation ist
\begin{align}
G( - \Omega) &= \left| G(\Omega) \right|  e^{- j  \varphi(\Omega)} \label{eq:8.25}
\end{align}
und \(y[k]\) kann vereinfacht werden zu
\begin{align}
{y\left[k\right]} &=  {\frac{A}{2}  \left|G\left(\Omega \right)\right| \left(e^{j \Omega  k}  e^{j \varphi \left(\Omega \right)} +e^{-j \Omega  k}  e^{-j \varphi \left(\Omega \right)} \right) \sigma \left[k\right]} \\
{} & {} & {-\frac{A}{2}  \left(\sum _{\kappa =k+1}^{\infty }g\left[\kappa \right]  e^{j \Omega  \left(k-\kappa \right)} +\sum _{\kappa =k+1}^{\infty }g\left[\kappa \right]  e^{-j \Omega  \left(k-\kappa \right)} \right) \sigma \left[k\right]} \\
{} &=  {A \left|G\left(\Omega \right)\right| \cos \left(\Omega  k+\varphi \left(\Omega \right)\right) \sigma \left[k\right]-A \left(\sum _{n=k+1}^{\infty }g\left[n\right]  \cos \left(\Omega  \left(k-n\right)\right)\right) \sigma \left[k\right]}
 \label{eq:8.26}
\end{align}
Der erste Term beschreibt die Systemreaktion auf die harmonische Anregung. Sie besitzt die gleiche Frequenz wie das Eingangssignal. Die Amplitude wird mit dem Betrag des Frequenzgangs \textbar \(G(Ω)\)\textbar{} multipliziert, die Phase ändert sich gegenüber dem Eingangssignal um die Phase des Frequenzgangs ϕ(Ω). Der zweite Term beschreibt das Einschwingen des Systems. Bei stabilen Systemen konvergiert die Impulsantwort \(g[k]\) für k → ∞ gegen null, sodass der zweite Summand bei stabilen Systemen gegen null geht. Damit gelten für zeitdiskrete Systeme bezüglich der Reaktion des Systems auf kausale, harmonische Anregungen dieselben Gesetzmäßigkeiten wie bei zeitkontinuierlichen Systemen.

\section{Grafische Darstellung des Frequenzgangs}\label{sec:8.2}

An dem Frequenzgang kann das Amplitudenverhältnis von Ein- und Ausgangssignal sowie die Phasenverschiebung als Funktion der normierten Kreisfrequenz Ω abgelesen werden. Die Beschreibung des Frequenzgangs erfolgt über den Amplitudengang und den Phasengang von Systemen. Sie werden in diesem Abschnitt berechnet und interpretiert. Der Frequenzgang zeitdiskreter Systeme entspricht der Fourier-Transformierten der Impulsantwort.
\begin{align}
G(\Omega) &= \sum_{k = - \infty}^{\infty}{g[ k]  e^{- j  k  \Omega}} \label{eq:8.27}
\end{align}
Wegen der Periodizität der Exponentialfunktion mit imaginärem Argument ist jeder Summand der Summe in Gleichung (\ref{eq:8.27}) periodisch in 2∙π. Damit ist der Frequenzgang eines zeitdiskreten Systems analog zu dem Spektrum von Signalfolgen periodisch in 2∙π. Deshalb ist es ausreichend, den Frequenzgang in einem Bereich von - π \ldots{} + π darzustellen. Der Frequenzgang kann in kartesischen Koordinaten oder Polar-Koordinaten dargestellt werden. Bei der Darstellung in kartesischen Koordinaten wird der Frequenzgang in Realteil \(G_{R}(Ω)\) und Imaginärteil \(G_{I}(Ω)\) zerlegt. Bei der Darstellung in Polarkoordinaten wird der Frequenzgang mit Betrag \textbar \(G(Ω)\)\textbar{} und Phase ϕ(Ω) dargestellt.
\begin{align}
G(\Omega) &= G_{R}(\Omega) + j  G_{I}(\Omega) = \left| G(\Omega) \right|  e^{j  \varphi(\Omega)} = A(\Omega)  e^{j  \varphi(\Omega)} \label{eq:8.28}
\end{align}
Der Betrag der Übertragungsfunktion \(A(Ω)\) wird als Amplitudengang und die Phase der Übertragungsfunktion ϕ(Ω) wird als Phasengang bezeichnet. Analog zu den Darstellungen des Frequenzgangs zeitkontinuierlicher Signale kann der Frequenzgang von Signalfolgen entweder als Ortskurve oder als Frequenzgangskennlinie dargestellt werden.

\subsection{Ortskurven}\label{sec:ortskurven}

Bei der Ortskurve wird der Frequenzgang in der komplexen Ebene abgebildet. Dazu werden für unterschiedliche Frequenzen - π ≤ Ω ≤ π Real- und Imaginärteil des Frequenzgangs bestimmt und die entsprechenden Punkte in die komplexe Ebene eingezeichnet.

\begin{example}{Rekursives Tiefpass-Filter}
\label{beispiel-rekursives-tiefpass-filter-3}

In Kapitel \ref{sec:8.1} wird der Frequenzgang des rekursiven Tiefpasses erster Ordnung mit 0 \textless{} GF \textless{} 1 berechnet.
\begin{align}
G(\Omega) &= \frac{1 - GF}{1 - GF  e^{- j  \Omega}} \label{eq:8.29}
\end{align}
Der Frequenzgang kann durch eine konjugiert komplexe Erweiterung in Real- und Imaginärteil aufgeteilt werden.
\begin{align}
{G\left(\Omega \right)} &=  {\frac{1-GF}{1-GF e^{-j \Omega } } =\frac{1-GF}{1-GF \cos \left(\Omega \right)+j GF \sin \left(\Omega \right)} } \\
{} &=  {\left(1-GF\right) \frac{1-GF \cos \left(\Omega \right)-j GF \sin \left(\Omega \right)}{\left(1-GF \cos \left(\Omega \right)\right)^{2} +GF^{2}  \sin ^{2} \left(\Omega \right)} } \\
{} &=  {\left(1-GF\right) \frac{1-GF \cos \left(\Omega \right)}{1+GF^{2} -2 GF \cos \left(\Omega \right)} -j \left(1-GF\right) \frac{GF \sin \left(\Omega \right)}{1+GF^{2} -2 GF \cos \left(\Omega \right)} }
 \label{eq:8.30}
\end{align}
Alternativ kann der Frequenzgang in Polarkoordinaten dargestellt werden, es ergibt sich der Betrag
\begin{align}
A(\Omega) &= \left| G(\Omega) \right| = \left| \frac{1 - GF}{1 - GF  \cos(\Omega) + j  GF  \sin(\Omega)} \right| = \frac{1 - GF}{\sqrt{1 + GF^{2} - 2  GF  \cos(\Omega)}} \label{eq:8.31}
\end{align}
und die Phase
\begin{align}
{\varphi \left(\Omega \right)} &=  {\varphi _{Z} \left(\Omega \right)-\varphi _{N} \left(\Omega \right)=\arctan \left(\frac{Im(Z)}{Re(Z)} \right)-\arctan \left(\frac{Im(N)}{Re(N)} \right)} \\
{} &=  {\arctan \left(\frac{0}{1-GF} \right)-\arctan \left(\frac{GF \sin \left(\Omega \right)}{1-GF \cos \left(\Omega \right)} \right)=-\arctan \left(\frac{GF \sin \left(\Omega \right)}{1-GF \cos \left(\Omega \right)} \right)}
 \label{eq:8.32}
\end{align}
Für den Frequenzbereich - π ≤ Ω ≤ π ergeben sich für GF = 0.5 und GF = 0.9 folgende Ortskurven.

\manualFigure
  [alt={Ortskurve eines rekursiven Tiefpassfilters mit GF = 0.5 und GF = 0.9}]
  {./media/image175.png}
  {Ortskurve eines rekursiven Tiefpassfilters mit GF = 0.5 und GF = 0.9}
  {}


Für das Filter mit GF = 0.5 wird die Ortskurve diskutiert. An den Stellen Ω = ± π und \(Ω = 0\) wird der Imaginärteil zu null, der Frequenzgang ist an diesen Stellen reell. Für die normierte Kreisfrequenz\\
Ω = ± π wird der Betrag minimal \(A = 1\)/√7 und für \(Ω = 0\) maximal (\(A = 1\)). Das Verhalten entspricht damit einem Tiefpass-Verhalten.

\end{example}

\subsection{Frequenzgangskennlinien}\label{sec:frequenzgangskennlinien}

Ortskurven stellen Real- und Imaginärteil in einem Schaubild dar, allerdings geht die Frequenzinformation bei der Darstellung verloren. Aus diesem Grund ist ihr Frequenzverhalten schwer zu interpretieren und wenig übersichtlich. Alternativ zur Ortskurve wird deshalb der Frequenzgang mit Frequenzgangskennlinien dargestellt. Dabei werden Betrag und Phase in separaten Diagrammen als Funktion der normierten Kreisfrequenz - π ≤ Ω ≤ π abgebildet. Frequenzgangskennlinien haben den Vorteil, dass der Zusammenhang zwischen Betrag beziehungsweise Phase und Frequenz direkt abgelesen werden kann.

\begin{example}{Rekursives Tiefpass-Filter}
\label{beispiel-rekursives-tiefpass-filter-4}

Wieder wird das Beispiel des rekursiven Tiefpasses mit dem Frequenzgang
\begin{align}
G(\Omega) &= \frac{1 - GF}{1 - GF  e^{- j  \Omega}} \label{eq:8.33}
\end{align}
aufgegriffen. Für den Frequenzbereich - π ≤ Ω ≤ π ergeben sich für GF = 0.5 und GF = 0.9 folgende Frequenzgangskennlinien.

\manualFigure
  [alt={Frequenzgangskennlinien eines rekursiven Tiefpassfilters mit GF = 0.5 und GF = 0.9}]
  {./media/image176.png}
  {Frequenzgangskennlinien eines rekursiven Tiefpassfilters mit GF = 0.5 und GF = 0.9}
  {}


Der Amplitudengang beider Filter weist bei der Frequenz Ω~=~0~mit A(0) = 1~sein Maximum auf. Er fällt bis zu den Stellen Ω = ± π ab. An allen drei Stellen ist die Phase ϕ = 0, weil die Übertragungsfunktion \(G(Ω)\) einen positiven reellen Wert aufweist. Bei dieser Darstellung ist die Tiefpass-Charakteristik direkt ablesbar. Es zeigt sich, dass die Grenzfrequenz und die Dämpfung im Sperrbereich von dem Gedächtnisfaktor GF abhängen.

Die normierte Kreisfrequenz \(Ω = 0\) entspricht dabei der Kreisfrequenz \(ω = 0\) und die normierten Kreisfrequenzen Ω = ± π entsprechen den Kreisfrequenzen ±~\(ω_{A}\)/2. Die Übertragungsfunktion des Systems hat einen Pol an der Stelle z = GF ≠ 0. Bei diesem Beispiel handelt es sich demnach um ein System mit unendlich langer Impulsantwort. Es wird als Inifinite-Impulse-Response (IIR) System bezeichnet. Der Frequenzgang von IIR-Systemen höherer Ordnung mit rationaler Übertragungsfunktion wird in Abschnitt \ref{sec:8.3} weiter vertieft.

\end{example}

\begin{example}{Gleitendes Mittelwertfilter}
\label{beispiel-gleitendes-mittelwertfilter}

Als zweites Beispiel wird der gleitende Mittelwert einer Signalfolge betrachtet. Es besitzt die Differenzengleichung
\begin{align}
y[ k] &= \frac{1}{5}  \left( u[ k] + u[ k - 1] + u[ k - 2] + u[ k - 3] + u[ k - 4] \right) \label{eq:8.34}
\end{align}
Das System hat eine endliche Impulsantwort \(g[k]\), die in Bild \ref{fig:4.2} dargestellt ist. Das Filter wird wegen der endlichen Impulsantwort als Finite-Impuls-Response (FIR) System bezeichnet. Mit der Verschiebungsregel der Fourier-Transformation ergibt sich die Darstellung im Frequenzbereich
\begin{align}
Y(\Omega) &= \frac{1}{5}  \left( 1 + e^{- j  \Omega} + e^{- 2  j  \Omega} + e^{- 3  j  \Omega} + e^{- 4  j  \Omega} \right)  U(\Omega) \label{eq:8.35}
\end{align}
Sie führt zu dem Frequenzgang des gleitenden Mittelwertes.
\begin{align}
{G\left(\Omega \right)} &=  {\frac{Y\left(\Omega \right)}{U\left(\Omega \right)} =\frac{1}{5}  \left(1+e^{-j \Omega } +e^{-2 j \Omega } +e^{-3 j \Omega } +e^{-4 j \Omega } \right)} \\
{} &=  {\frac{1}{5}  e^{-2 j \Omega }  \left(e^{2 j \Omega } +e^{j \Omega } +1+e^{-j \Omega } +e^{-2 j \Omega } \right)} \\
{} &=  {\frac{1}{5}  e^{-2 j \Omega }  \left(1+2 \cos \left(\Omega \right)+2 \cos \left(2 \Omega \right)\right)} \\
{} &=  {\frac{1}{5}  \left(1+2 \cos \left(\Omega \right)+2 \cos \left(2 \Omega \right)\right) \left(\cos \left(2 \Omega \right)-j \sin \left(2 \Omega \right)\right)}
 \label{eq:8.36}
\end{align}
Der Frequenzgang ist in Bild \ref{fig:8.3} als Frequenzgangskennlinie dargestellt.

\manualFigure
  [alt={Frequenzgangskennlinie für den gleitenden Mittelwert der Ordnung \(N = 5\)}]
  {./media/image178.png}
  {Frequenzgangskennlinie für den gleitenden Mittelwert der Ordnung \(N = 5\)}
  {\label{fig:8.3}}


Der Amplitudengang zeigt, dass die Spektralanteile mit niedriger Frequenz das Filter ungedämpft passieren können, während Spektralanteile mit höherer Frequenz deutlich abgeschwächt werden. Es handelt sich demnach ebenfalls um ein Tiefpass-Filter.

Der Phasengang besteht aus Abschnitten, die jeweils über eine Gerade beschrieben werden können. Das Springen der Phase um π ergibt sich aus Vorzeichenwechseln des Frequenzgangs. Die Sprünge finden bei diesem Beispiel immer dann statt, wenn der Amplitudengang den Wert 0 aufweist, sodass für alle relevanten Frequenzen eine eindeutige Phasenbeziehung gegeben ist.

Auf die Linearität der Phase und die Auswirkung auf das Ausgangssignal wird im Zusammenhang mit FIR-Systemen in Abschnitt \ref{sec:8.4} näher eingegangen.

\end{example}

\subsection{Bode-Diagramme}\label{sec:bode-diagramme}

Im zeitkontinuierlichen Bereich werden Frequenzgangskennlinien als Bode-Diagramme dargestellt. Dabei wird die Frequenzachse logarithmisch skaliert und der Amplitudengang in Dezibel
\begin{align}
a(\Omega) &= 20  \log\left( A(\Omega) \right)\ dB \label{eq:8.37}
\end{align}
dargestellt. Diese Darstellungsform ist auch im zeitdiskreten Bereich gebräuchlich. Allerdings muss dabei berücksichtigt werden, dass der Frequenzgang periodisch in 2π ist. Außerdem kann er wegen des logarithmischen Maßstabs nur positive Frequenzen Ω dargestellt werden. Insgesamt beschränkt sich damit der Frequenzbereich auf 0 \textless{} Ω ≤ π.

Das Bode-Diagramm wird bevorzugt für einen Vergleich von analogen und digitalen Systemen eingesetzt. Dazu wird die normierte Frequenzachse Ω über
\begin{align}
\omega &= \frac{\Omega}{T_{A}} \label{eq:8.38}
\end{align}
in eine Frequenzachse ω überführt. Damit können die Frequenzgänge des zeitkontinuierlichen und des zeitdiskreten Systems in einem Diagramm verglichen werden.

\begin{example}{Vergleich eines zeitkontinuierlichen und zeitdiskreten Tiefpass-Filters}
\label{beispiel-vergleich-eines-zeitkontinuierlichen-und-zeitdiskreten-tiefpass-filters}

Ein rekursives zeitdiskretes Filter mit dem Frequenzgang
\begin{align}
G_{D}(\Omega) &= \frac{1 - GF}{1 - GF  e^{- j  \Omega}} \label{eq:8.39}
\end{align}
mit einem Gedächtnisfaktor GF = 0.9 wird mit einer Abtastzeit \(T_{A} = 1\,\text{ms}\) betrieben. Die Filterwirkung wird mit einem Tiefpass erster Ordnung verglichen, der eine Zeitkonstante \(T = 10\,\text{ms}\) aufweist.
\begin{align}
G_{K}(\omega) &= \frac{1}{1 + j  \omega  T} \label{eq:8.40}
\end{align}
Bild \ref{fig:8.4} stellt das Bode-Diagramm für die Tiefpass-Filter dar.

\manualFigure
  [alt={Bode-Diagramm eines zeitdiskreten rekursiven Tiefpassfilters mit GF = 0.9 und \(T_{A} = 1\,\text{ms}\)\\
und eines zeitkontinuierlichen Tiefpass-Filters mit \(T = 10\,\text{ms}\)}]
  {./media/image179.png}
  {Bode-Diagramm eines zeitdiskreten rekursiven Tiefpassfilters mit GF = 0.9 und \(T_{A} = 1\,\text{ms}\)\\
und eines zeitkontinuierlichen Tiefpass-Filters mit \(T = 10\,\text{ms}\)}
  {\label{fig:8.4}}


Die Amplitudengänge sind insbesondere in dem niedrigen Frequenzbereich identisch, erst im Bereich der Abtastfrequenz \(ω_{A}\)/2 ergeben sich signifikante Abweichungen. Die Phasengänge sind bereits in niedrigen Frequenzbereich unterschiedlich.

\end{example}

\section{Pol-Nullstellen-Diagramm und Frequenzgang eines Systems}\label{sec:8.3}

Die bislang diskutierten linearen, zeitinvarianten Systeme führen zu linearen Differenzengleichungen mit konstanten Koeffizienten. Ihr Systemverhalten wird mit der z-Transformation durch die Übertragungsfunktion
\begin{align}
G(z) &= \frac{B(z)}{A(z)} = \frac{\sum_{m = 0}^{M}{b_{m}  z^{m}}}{\sum_{n = 0}^{N}{a_{n}  z^{n}}} = \frac{b_{M}}{a_{N}}  \frac{\prod_{m = 1}^{M}\left( z - \beta_{m} \right)}{\prod_{n = 1}^{N}\left( z - \alpha_{n} \right)} \label{eq:8.41}
\end{align}
beschrieben. Bei stabilen Systemen liegen die Pole \(α_{n}\) der Übertragungsfunktion im Einheitskreis. In dem Fall ergibt sich der Frequenzgang zu
\begin{align}
G(\Omega) &= \left. \ G(z) \right|_{z = e^{j  \Omega}} = A(\Omega)  e^{j  \varphi(\Omega)} = \frac{b_{M}}{a_{N}}  \frac{\prod_{m = 1}^{M}\left( e^{j  \Omega} - \beta_{m} \right)}{\prod_{n = 1}^{N}\left( e^{j  \Omega} - \alpha_{n} \right)} \label{eq:8.42}
\end{align}
Der Amplitudengang des Systems kann logarithmisch in dB angegeben werden. Der logarithmische Amplitudengang \(a(Ω)\) ergibt sich mit den Rechenregeln des Logarithmus zu
\begin{align}
{a\left(\Omega \right)} &=  {20 \log \left(\left|\frac{b_{M} }{a_{N} }  \frac{\prod _{m=1}^{M}\left(e^{j \Omega } -\beta _{m} \right) }{\prod _{n=1}^{N}\left(e^{j \Omega } -\alpha _{n} \right) } \right|\right)} \\
{} &=  {20 \log \left(\left|\frac{b_{M} }{a_{N} } \right|\right)+\sum _{m=1}^{M}20 \log  \left(\left|e^{j \Omega } -\beta _{m} \right|\right)-\sum _{n=1}^{N}20 \log  \left(\left|e^{j \Omega } -\alpha _{n} \right|\right)}
 \label{eq:8.43}
\end{align}
Abgesehen von dem ersten konstanten Term setzt sich der logarithmische Amplitudengang aus Summanden zusammen, die immer dieselbe Struktur aufweisen. Im Zähler stehende Linearfaktoren gehen positiv in den logarithmischen Amplitudengang ein. Linearfaktoren, die im Nenner stehen, gehen negativ in den logarithmischen Amplitudengang ein. Auch die Phase kann nach den Rechenregeln komplexer Exponentialfunktionen über eine Summe ausgedrückt werden.
\begin{align}
\varphi(\Omega) &= \arg\left( \frac{b_{M}}{a_{N}} \right) + \sum_{m = 1}^{M}{\arg\left( e^{j  \Omega} - \beta_{m} \right)} - \sum_{n = 1}^{N}{\arg\left( e^{j  \Omega} - \alpha_{n} \right)} \label{eq:8.44}
\end{align}
Im Zähler stehende Linearfaktoren liefern positive und im Nenner stehende Linearfaktoren negative Beiträge. Aus diesem Grund ist die Interpretation des Frequenzgangs eines Linearfaktors der Form
\begin{align}
G(\Omega) &= e^{j  \Omega} - \beta \label{eq:8.45}
\end{align}
von Interesse. Die Frequenzgänge aller stabilen LTI-Systeme können mit dieser Grundlage erstellt werden.
\begin{application}{Komplexe Exponentialfolge}
  Im Online-Portal \emph{Systemtheorie Online} verdeutlicht die Applikation \emph{Komplexe Exponentialfolge} den Zusammenhang zwischen der Pollage in der komplexen Ebene und dem Frequenzgang des zugehörigen Systems.
  \tcblower
  \url{h-ka.de/eit/digitale-signalverarbeitung/übungen/kapitel2}
\end{application}
\subsection{Frequenzgang eines Systems mit einer Nullstelle}\label{sec:frequenzgang-eines-systems-mit-einer-nullstelle}

Im Folgenden wird unabhängig von der Realisierbarkeit eine Nullstelle im Einheitskreis
\begin{align}
\beta &= r_{0}  e^{j  \Omega_{0}} \label{eq:8.46}
\end{align}
diskutiert. Mit dieser Definition ergibt sich der Frequenzgang
\begin{align}
{G\left(\Omega \right)} &=  {e^{j \Omega } -\beta =e^{j \Omega } -r_{0}  e^{j \Omega _{0} } =e^{j \Omega }  \left(1-r_{0}  e^{j \left(\Omega _{0} -\Omega \right)} \right)} \\
{} &=  {e^{j \Omega }  \left(1-r_{0}  \cos \left(\Omega _{0} -\Omega \right)-j r_{0}  \sin \left(\Omega _{0} -\Omega \right)\right)}
 \label{eq:8.47}
\end{align}
Der Frequenzgang hat einen Betrag von
\begin{align}
A(\Omega) &= \sqrt{\left( 1 - r_{0}  \cos\left( \Omega_{0} - \Omega \right) \right)^{2} + r_{0}^{2}  \sin^{2}\left( \Omega_{0} - \Omega \right)} = \sqrt{1 + r_{0}^{2} - 2  r_{0}  \cos\left( \Omega_{0} - \Omega \right)} \label{eq:8.48}
\end{align}
und eine Phase von
\begin{align}
\varphi(\Omega) &= \Omega + \arctan\left( \frac{- r_{0}  \sin\left( \Omega_{0} - \Omega \right)}{1 - r_{0}  \cos\left( \Omega_{0} - \Omega \right)} \right) = \Omega - \arctan\left( \frac{r_{0}  \sin\left( \Omega_{0} - \Omega \right)}{1 - r_{0}  \cos\left( \Omega_{0} - \Omega \right)} \right) \label{eq:8.49}
\end{align}
Beide Größen können in der komplexen Ebene geometrisch interpretiert werden. Bild \ref{fig:8.5} verdeutlicht die Geometrie für eine feste Kreisfrequenz Ω.

\begin{figure}
  \centering
  \includegraphics[alt={Darstellung von Betrag und Phase im Pol-Nullstellen-Diagramm\\
für ein System mit einer komplexen Nullstelle im Einheitskreis}]{./media/image182.png}
  \caption{Darstellung von Betrag und Phase im Pol-Nullstellen-Diagramm\\
für ein System mit einer komplexen Nullstelle im Einheitskreis}
  \label{fig:8.5}
\end{figure}

Der Betrag der Übertagungsfunktion \(A(Ω)\) ergibt sich nach dem Satz des Pythagoras aus dem rechtwinkligen Dreieck mit den Kantenlängen \textbar1 - r\textsubscript{0}cos(Ω - \(Ω_{0}\))\textbar{} und \textbar r\textsubscript{0}sin(Ω - \(Ω_{0}\))\textbar. Da der Einheitskreis periodisch durchlaufen wird, ergibt sich ein periodischer Amplitudengang. Die Phase des Frequenzgangs ϕ(Ω) kann ebenfalls über das rechtwinklige Dreieck berechnet werden. Um den Einfluss der Lage der Nullstelle zu verdeutlichen, wird das Systemverhalten anhand zweier Parameterstudien diskutiert.

\subsection{Variation des Abstandes der Nullstelle vom Koordinatenursprung}\label{sec:8.3.2}

In der ersten Parameterstudie wird der Abstand der Nullstelle vom Koordinatenursprung in den Schritten \(r_{1} = 0.5\), \(r_{2} = 0.75\) und \(r_{3} = 1\) variiert und der Einfluss auf Amplitudengang und Phasengang diskutiert. Der Winkel \(Ω_{0}\), den die Nullstelle in der komplexen Ebene besitzt, bleibt konstant auf\\
\(Ω_{0} = π\)/2. Zur besseren Übersicht stellt Bild \ref{fig:8.6} die Lage der Nullstelle bei Variation des Abstandes vom Koordinatenursprung in der komplexen Ebene dar.

\begin{figure}
  \centering
  \includegraphics[alt={Lage der Nullstelle bei Variation des Abstandes vom Koordinatenursprung, \(Ω_{0} = π\)/2}]{./media/9_1_lage_der_nullstelle_bei_variatio.png}
  \caption{Lage der Nullstelle bei Variation des Abstandes vom Koordinatenursprung, \(Ω_{0} = π\)/2}
  \label{fig:8.6}
\end{figure}

Der Amplitudengang und der Phasengang berechnen sich für \(Ω_{0} = π\)/2 zu
\begin{align}
A_{n}(\Omega) &= \sqrt{1 + r_{n}^{2} - 2  r_{n}  \cos\left( \frac{\pi}{2} - \Omega \right)} \label{eq:8.50}
\end{align}
beziehungsweise zu
\begin{align}
\varphi_{n}(\Omega) &= \Omega - \arctan\left( \frac{r_{n}  \sin\left( \frac{\pi}{2} - \Omega \right)}{1 - r_{n}  \cos\left( \frac{\pi}{2} - \Omega \right)} \right) \label{eq:8.51}
\end{align}
Für die verschiedenen Lagen der Nullstelle ergeben sich die in Bild \ref{fig:8.7} dargestellten Amplituden- und Phasengänge.

\begin{figure}
  \centering
  \includegraphics[alt={Amplituden- und Phasengang bei variiertem Abstand der Nullstelle vom Koordinatenursprung, \(Ω_{0} = π\)/2}]{./media/9_2_amplituden_und_phasengang_bei_va.png}
  \caption{Amplituden- und Phasengang bei variiertem Abstand der Nullstelle vom Koordinatenursprung, \(Ω_{0} = π\)/2}
  \label{fig:8.7}
\end{figure}

Der Amplitudengang weist sein Minimum bei der Frequenz \(Ω = Ω_{0}\) = π/2 auf. In dem Fall wird das Argument der Kosinus-Funktion zu null und damit der Amplitudengang minimal. Der minimale Wert des Amplitudengangs beträgt
\begin{align}
A_{MIN}(\Omega) &= \sqrt{1 + r_{n}^{2} - 2  r_{n}} = 1 - r_{n} \label{eq:8.52}
\end{align}
Je näher die Nullstelle am Einheitskreis liegt, desto kleiner wird der Amplitudengang an der Stelle\\
\(Ω = Ω_{0}\). Bei dem Radius \(r_{3} = 1\) wird der Amplitudengang zu null. Das Maximum des Amplitudengangs wird bei der Frequenz \(Ω = Ω_{0}\) - π= - π/2 erreicht. Bei dieser Frequenz wird das Argument der Kosinus-Funktion zu π und damit der Amplitudengang maximal. Der maximale Wert des Amplitudengangs beträgt
\begin{align}
A_{MAX}(\Omega) &= \sqrt{1 + r_{n}^{2} + 2  r_{n}} = 1 + r_{n} \label{eq:8.53}
\end{align}
Der Phasengang springt für den Radius \(r_{3} = 1\) an der Stelle \(Ω = Ω_{0}\) = π/2, da das Argument der Arkustangensfunktion unendlich wird. Der Phasengang ist für \(r_{3}\) damit unstetig, ist aber ansonsten linear. Je weiter die Nullstelle vom Einheitskreis entfernt ist, desto stärker wird dieser Sprung abgeflacht. Unabhängig vom Radius beträgt der Phasengang bei der Frequenz \(Ω = Ω_{0}\) = - π/2
\begin{align}
\varphi\left( - \frac{\pi}{2} \right) &= - \frac{\pi}{2} - \arctan\left( \frac{0}{1 + r} \right) = - \frac{\pi}{2} \label{eq:8.54}
\end{align}
\subsection{Variation des Phasenwinkels der Nullstelle}\label{sec:variation-des-phasenwinkels-der-nullstelle}

In der zweiten Parameterstudie wird die Phasenlage der Nullstelle mit \(Ω_{1} = 0\), \(Ω_{2} = π\)/2 und \(Ω_{3} = π\) variiert. Der Radius bleibt konstant auf dem Wert \(r_{0} = 0.7\). Zur besseren Übersicht stellt Bild \ref{fig:8.8} die Lage der Nullstelle bei Variation des Phasenwinkes \(Ω_{n}\) der Nullstelle dar.

\includegraphics[width=3.14961in,height=2.24803in,alt={Ein Bild, das Diagramm, Text, Reihe, Screenshot enthält. Automatisch generierte Beschreibung}]{./media/image185.png}

\protect\phantomsection\label{fig:8.8}{}Bild .: Lage der Nullstelle bei Variation des Phasenwinkes der Nullstelle, \(r_{0} = 0.7\)

Der Amplitudengang und der Phasengang berechnen sich für \(r_{0} = 0.7\) zu
\begin{align}
A_{n}(\Omega) &= \sqrt{1 + 0.7^{2} - 2  0.7  \cos\left( \Omega_{n} - \Omega \right)} = \sqrt{1.49 - 1.4  \cos\left( \Omega_{n} - \Omega \right)} \label{eq:8.55}
\end{align}
beziehungsweise zu
\begin{align}
\varphi_{n}(\Omega) &= \Omega - \arctan\left( \frac{0.7  \sin\left( \Omega_{0} - \Omega \right)}{1 - 0.7  \cos\left( \Omega_{0} - \Omega \right)} \right) \label{eq:8.56}
\end{align}
Für die verschiedenen Lagen der Nullstelle ergeben sich die in Bild \ref{fig:8.9} dargestellten Amplituden- und Phasengänge.

\begin{figure}
  \centering
  \includegraphics[alt={Amplituden- und Phasengang bei variiertem Phasenwinkel α der Nullstelle, \(r = 0.7\)}]{./media/9_3_amplituden_und_phasengang_bei_va.png}
  \caption{Amplituden- und Phasengang bei variiertem Phasenwinkel α der Nullstelle, \(r = 0.7\)}
  \label{fig:8.9}
\end{figure}

Durch die unterschiedlichen Phasenwinkel \(Ω_{0}\) verschiebt sich der Amplitudengang um \(Ω_{0}\) nach rechts, ansonsten bleibt Amplitudengang gleich. Der Phasengang verschiebt sich mit steigendem Phasenwinkel \(Ω_{0}\) einerseits nach rechts, aber wegen des Summanden Ω in Gleichung (\ref{eq:8.56}) auch nach oben.

\subsection{Interpretationsbeispiel}\label{sec:interpretationsbeispiel}

Die Berechnung und Interpretation des Frequenzgangs von Systemen höherer Ordnung wird an einem rekursiven IIR-System erläutert. Dabei wird außerdem der Zusammenhang von Zeitbereich, z-Bereich und Frequenzbereich vertieft.

Gegeben ist ein zeitdiskretes System, das mit einer Abtastzeit von \(T_{A} = 5\,\text{ms}\) arbeitet und die Übertragungsfunktion
\begin{align}
G(z) &= \frac{z^{2} - 0.81}{z^{2} + 0.49} = \frac{(z - 0.9)  (z + 0.9)}{(z - j  0.7)  (z + j  0.7)} \label{eq:8.57}
\end{align}
besitzt. Das System hat damit zwei Pole und zwei Nullstellen. Da der Zählergrad genauso groß ist wie der Nennergrad, muss für die Berechnung der Impulsantwort zunächst eine Polynomdivision durchgeführt werden. Eine anschließende Partialbruchzerlegung führt zu
\begin{align}
G(z) &= 1 + \frac{j  \text{0.9286}}{z - j  0.7} - \frac{j  \text{0.9286}}{z + j  0.7} = 1 + j  \text{0.9286}  \left( \frac{z}{z - j  0.7} - \frac{z}{z + j  0.7} \right)  z^{- 1} \label{eq:8.58}
\end{align}
Damit lautet die Impulsantwort des Systems
\begin{align}
{g\left[k\right]} &=  {\delta \left[k\right]+j 0.9286 \left(\left(j 0.7\right)^{k-1} -\left(-j 0.7\right)^{k-1} \right) \sigma \left[k-1\right]} \\
{} &=  {\delta \left[k\right]+j 0.9286 0.7^{k-1}  \left(e^{j \frac{\pi }{2}  \left(k-1\right)} -e^{-j \frac{\pi }{2}  \left(k-1\right)} \right) \sigma \left[k-1\right]} \\
{} &=  {\delta \left[k\right]+j 0.9286 0.7^{k-1}  2 j \sin \left(\frac{\pi }{2}  \left(k-1\right)\right) \sigma \left[k-1\right]} \\
{} &=  {\delta \left[k\right]-0.9286 0.7^{k-1}  2 \sin \left(\frac{\pi }{2}  \left(k-1\right)\right) \sigma \left[k-1\right]}
 \label{eq:8.59}
\end{align}
Da der Zählergrad genauso groß ist wie der Nennergrad, ist das System sprungfähig. Die Pol-Nullstellenverteilung sowie die Impulsantwort sind in Bild \ref{fig:8.10} dargestellt.

\begin{figure}
  \centering
  \includegraphics[alt={Pol-Nullstellenverteilung sowie die Impulsantwort des IIR-Systems 2. Ordnung}]{./media/9_4_polnullstellenverteilung_sowie_d.png}
  \caption{Pol-Nullstellenverteilung sowie die Impulsantwort des IIR-Systems 2. Ordnung}
  \label{fig:8.10}
\end{figure}

Die Pole des Systems haben einen Betrag von 0.7. Damit ist das System stabil, und der Frequenzgang kann durch die Substitution z = \(e^{jΩ}\) bestimmt werden.
\begin{align}
G(\Omega) &= \frac{\left( 1 - 0.9  e^{- j  \Omega} \right)  \left( 1 + 0.9  e^{- j  \Omega} \right)}{\left( 1 - j  0.7  e^{- j  \Omega} \right)  \left( 1 + j  0.7  e^{- j  \Omega} \right)} \label{eq:8.60}
\end{align}
Der Frequenzgang setzt sich aus vier Linearfaktoren zusammen. Nach den Darstellungen zum Frequenzgang ergibt sich der logarithmische Amplitudengang zu
\begin{align}
{a\left(\Omega \right)} &=  {10 \log \left(1.81-1.8 \cos \left(\Omega \right)\right)+10 \log \left(1.81+1.8 \cos \left(\Omega \right)\right)} \\
{} & {} & {-10 \log \left(1.49+1.4 \sin \left(\Omega \right)\right)-10 \log \left(1.49-1.4 \sin \left(\Omega \right)\right)}
 \label{eq:8.61}
\end{align}
und der Phasengang zu
\begin{align}
{\varphi \left(\Omega \right)} &=  {\arctan \left(\frac{0.9 \sin \left(\Omega \right)}{1-0.9 \cos \left(\Omega \right)} \right)+\arctan \left(\frac{-0.9 \sin \left(\Omega \right)}{1+0.9 \cos \left(\Omega \right)} \right)} \\
{} & {} & {-\arctan \left(\frac{-0.7 \cos \left(\Omega \right)}{1-0.7 \sin \left(\Omega \right)} \right)-\arctan \left(\frac{0.7 \cos \left(\Omega \right)}{1-0.7 \sin \left(\Omega \right)} \right)}
 \label{eq:8.62}
\end{align}
Bild \ref{fig:8.11} stellt den Amplituden- und Phasengang des Systems dar.

\begin{figure}
  \centering
  \includegraphics[alt={Frequenzgang des IIR-Systems 2. Ordnung}]{./media/9_5_frequenzgang_des_iirsystems_2_or.png}
  \caption{Frequenzgang des IIR-Systems 2. Ordnung}
  \label{fig:8.11}
\end{figure}

Bei dem System handelt es sich um einen zeitdiskret realisierten Bandpass.

\section{\texorpdfstring{Interpretation des Phasengangs eines Systems }{Interpretation des Phasengangs eines Systems }}\label{sec:8.4}

Der Frequenzgang eines Systems setzt sich aus Amplitudengang und Phasengang zusammen. Während der Amplitudengang anschaulich interpretiert werden kann, erscheint die Interpretation des Phasengangs kompliziert. In diesem Abschnitt wird deshalb der Phasengang von Systemen ausführlich diskutiert. Der Begriff der Phase wird im Zusammenhang mit Zeigerdiagrammen und harmonischen Anregungen eingeführt. Bei der harmonischen Anregung eines Systems entspricht die Phasenverschiebung einer zeitlichen Verschiebung des Signals. Der im Zusammenhang mit dem Frequenzgang berechnete Phasengang ergibt sich aus der Fourier-Transformation. Es wird an Beispielen gezeigt, dass mit dem Phasengang eines Systems die Verzögerung eines Eingangssignals verbunden ist, aus der der Begriff der sogenannten Gruppenlaufzeit resultiert. Einen Sonderfall nehmen Systeme mit linearem Phasengang ein. Sie besitzen eine konstante Gruppenlaufzeit.

\subsection{Phasengang eines Totzeitgliedes}\label{sec:phasengang-eines-totzeitgliedes}

Zur Verdeutlichung des Begriffes der Phase wird ein ideales Totzeitglied hinsichtlich seines Phasengangs analysiert. Unter der Annahme, dass die Totzeit ein Vielfaches der Abtastzeit \(T_{A}\) ist, hat das ideale Totzeitglied die Impulsantwort
\begin{align}
g[ k] &= \delta\left[ k - k_{0} \right] \label{eq:8.63}
\end{align}
Eine am Eingang angelegte Signalfolge bleibt in ihrer Form vollständig erhalten, sie wird lediglich um \(k_{0}\) nach rechts verschoben. Der Frequenzgang des idealen Totzeitglieds ergibt sich aus
\begin{align}
G(\Omega) &= e^{- j  \Omega  k_{0}} \label{eq:8.64}
\end{align}
Er hat den Betrag von 1 und eine Phase von
\begin{align}
\varphi(\Omega) &= - \Omega  k_{0} \label{eq:8.65}
\end{align}
Die Phase ϕ ist eine lineare Funktion der normierten Kreisfrequenz Ω. Bild \ref{fig:8.12} zeigt den Phasengang und die Sprungantwort des idealen Totzeitgliedes.

\begin{figure}
  \centering
  \includegraphics[alt={Phasengang und Sprungantwort eines idealen Totzeitglieds mit \(k_{0} = 5\)}]{./media/9_6_phasengang_und_sprungantwort_ein.png}
  \caption{Phasengang und Sprungantwort eines idealen Totzeitglieds mit \(k_{0} = 5\)}
  \label{fig:8.12}
\end{figure}

Die Sprungantwort zeigt, dass sich die Signalform bei einem idealen Totzeitglied mit linearem Phasengang nicht ändert.

\subsection{Anregung eines Systems mit harmonischen Signalen}\label{sec:anregung-eines-systems-mit-harmonischen-signalen}

Ein LTI-System reagiert im Fall einer harmonischen Anregung mit einem Ausgangssignal gleicher Frequenz, aber gegebenenfalls unterschiedlicher Amplitude und Phase. Bild \ref{fig:8.13} zeigt ein Beispiel für Ein- und Ausgangsfolge bei harmonischer Anregung.

\begin{figure}
  \centering
  \includegraphics[alt={Ein- und Ausgangsfolge eines Systems bei harmonischer Anregung}]{./media/9_7_ein_und_ausgangsfolge_eines_syst.png}
  \caption{Ein- und Ausgangsfolge eines Systems bei harmonischer Anregung}
  \label{fig:8.13}
\end{figure}

Die beiden Signale können mit den Folgen
\begin{align}
u[ k] &= 5  \cos\left( \frac{2  \pi}{10}  k \right) = 5  \cos\left( \Omega_{0}  k \right) \label{eq:8.66}
\end{align}
und
\begin{align}
{y[ k] = 4  \cos\left( \frac{2  \pi}{10}  (k - 2) \right) = 4  \cos\left( \frac{2  \pi}{10}  k - \frac{4  \pi}{10} \right)
}{= 4  \cos\left( \Omega_{0}  k + \varphi \right) = 4  \cos\left( \Omega_{0}  \left( k - k_{0} \right) \right)} \label{eq:8.67}
\end{align}
beschrieben werden. Die Phase ϕ des Ausgangssignals kann über eine Verschiebung \(k_{0}\) und der Frequenz \(Ω_{0}\) dargestellt werden als
\begin{align}
\varphi(\Omega) &= - \Omega_{0}  k_{0} \label{eq:8.68}
\end{align}
Bei der Anregung mit einer harmonischen Schwingung der Frequenz \(Ω_{0}\) ist für die Berechnung des Ausgangssignals nur der Frequenzgang des Systems an dieser Stelle \(Ω_{0}\) von Bedeutung. Deshalb hat eine Umrechnung der Zeitverschiebung \(k_{0}\) mit der Frequenz im Fall einer harmonischen Anregung lediglich normierenden Charakter. Die Phase ist in dem Fall gleichbedeutend mit einer zeitlichen Verschiebung der harmonischen Funktion.

\subsection{\texorpdfstring{Anregung eines Systems mit zwei harmonischen Signalen unterschiedlicher Frequenz }{Anregung eines Systems mit zwei harmonischen Signalen unterschiedlicher Frequenz }}\label{anregung-eines-systems-mit-zwei-harmonischen-signalen-unterschiedlicher-frequenz}

Die Vorstellung der Phase als Maß für eine zeitliche Verschiebung wird am Beispiel zweier harmonischer Schwingungen weiter ausgebaut. Gegeben ist ein Signal, das aus der Summe zweier Folgen
\begin{align}
u_{1}[ k] &= U_{1}  \sin\left( \Omega_{1}  k \right) \label{eq:8.69}
\end{align}
und
\begin{align}
u_{2}[ k] &= U_{2}  \sin\left( \Omega_{2}  k \right) \label{eq:8.70}
\end{align}
besteht. Das Signal durchläuft ein System mit dem Frequenzgang \(G(Ω)\). Das Ausgangssignal ergibt sich aus der Summe der Folgen
\begin{align}
y_{1}[ k] &= U_{1}  A\left( \Omega_{1} \right)  \sin\left( \Omega_{1}  k + \varphi\left( \Omega_{1} \right) \right) \label{eq:8.71}
\end{align}
und
\begin{align}
y_{2}[ k] &= U_{2}  A\left( \Omega_{2} \right)  \sin\left( \Omega_{2}  k + \varphi\left( \Omega_{2} \right) \right) \label{eq:8.72}
\end{align}
Damit ein Signal dieselbe Form behält, müssen die Signalanteile unterschiedlicher Frequenz um die gleiche Zeit \(k_{0}\) verzögert werden. Unter dieser Bedingung ergibt sich für die beiden Signale
\begin{align}
y_{1}[ k] &= U_{1}  A\left( \Omega_{1} \right)\sin\left( \Omega_{1}  k + \varphi\left( \Omega_{1} \right) \right) = U_{1}  A\left( \Omega_{1} \right)  \sin\left( \Omega_{1}  \left( k - k_{0} \right) \right) \label{eq:8.73} \\
y_{2}[ k] &= X_{2}  A\left( \Omega_{2} \right)  \sin\left( \Omega_{2}  k + \varphi\left( \Omega_{2} \right) \right) = X_{2}  A\left( \Omega_{2} \right)  \sin\left( \Omega_{2}  \left( k - k_{0} \right) \right) \label{eq:8.74}
\end{align}
und entsprechend für die beiden Phasen
\begin{align}
\varphi\left( \Omega_{1} \right) &= - k_{0}  \Omega_{1} \label{eq:8.75} \\
\varphi\left( \Omega_{2} \right) &= - k_{0}  \Omega_{2} \label{eq:8.76}
\end{align}
Dieser Gedankengang zeigt, dass ein System mit linearer Phase Signale nur um eine konstante Zeit \(k_{0}\) verschiebt. Eine Phasenverzerrung des Signals, bei dem unterschiedliche Spektralbereiche zeitlich gesehen unterschiedlich weit verschoben werden, findet nicht statt.

\subsection{Gruppenlaufzeit eines Systems}\label{sec:8.4.4}

Die vorangegangenen Beispiele zeigen, dass der Phasengang für eine Verzögerung der Signale verantwortlich ist. Diese Verzögerung wird mit der sogenannten Gruppenlaufzeit beschrieben. Aus physikalischen Gründen ist sie als die normierte Ableitung des Phasengangs definiert.
\begin{align}
T_{G}(\Omega) &= - \frac{d\varphi}{d\Omega}  T_{A} \label{eq:8.77}
\end{align}
Im Allgemeinen ist die Gruppenlaufzeit eines Systems eine Funktion der Frequenz Ω und damit nicht konstant. Ist der Phasengang linear, ist seine Ableitung konstant. Damit ist auch die Gruppenlaufzeit eines Systems mit linearer Phase konstant.

\begin{example}{Rekursives Tiefpass-Filter}
\label{beispiel-rekursives-tiefpass-filter-5}

In Kapitel \ref{sec:8.1} wird der Frequenzgang des rekursiven Tiefpasses erster Ordnung berechnet.
\begin{align}
G(\Omega) &= \frac{1 - GF}{1 - GF  e^{- j  \Omega}} \label{eq:8.78}
\end{align}
Der Phasengang wird in Gleichung (\ref{eq:8.32}) berechnet zu
\begin{align}
\varphi(\Omega) &= - \arctan\left( \frac{GF  \sin(\Omega)}{1 - GF  \cos(\Omega)} \right) \label{eq:8.79}
\end{align}
Die Ableitung nach Ω führt zu der Gruppenlaufzeit von
\begin{align}
{T_{G} \left(\Omega \right)} &=  {-\frac{d\varphi }{d\Omega }  T_{A} =\frac{\frac{GF \cos \left(\Omega \right)}{1-GF \cos \left(\Omega \right)} -\frac{GF^{2}  \sin ^{2} \left(\Omega \right)}{\left(1-GF \cos \left(\Omega \right)\right)^{2} } }{1+\frac{GF^{2}  \sin ^{2} \left(\Omega \right)}{\left(1-GF \cos \left(\Omega \right)\right)^{2} } }  T_{A} } \\
{} &=  {\frac{GF \cos \left(\Omega \right)-GF^{2} }{1-2 GF \cos \left(\Omega \right)+GF^{2} }  T_{A} }
 \label{eq:8.80}
\end{align}
Bild \ref{fig:8.14} zeigt die Gruppenlaufzeit eines rekursiven Filters 1. Ordnung für GF = 0.9 und eine Abtastzeit von \(T_{A} = 1\,\text{ms}\).

\manualFigure
  [alt={Gruppenlaufzeit eines rekursiven Filters (GF = 0.9, \(T_{A} = 1\,\text{ms}\))}]
  {./media/image195.png}
  {Gruppenlaufzeit eines rekursiven Filters (GF = 0.9, \(T_{A} = 1\,\text{ms}\))}
  {\label{fig:8.14}}


Das Filter weist nur im Bereich kleiner Frequenzen ω eine konstante Gruppenlaufzeit auf, bei höheren Frequenzen ändert sich die Gruppenlaufzeit. Der in Bild \ref{fig:8.4} abgebildete Amplitudengang zeigt, dass es sich bei dem Filter um ein Tiefpassfilter handelt, das im niedrigen Frequenzbereich eine näherungsweise konstante Gruppenlaufzeit besitzt. Der Spektralbereich, der aufgrund der variablen Gruppenlaufzeit Phasenverzerrungen erfährt, wird gedämpft.

\end{example}

\begin{example}{Totzeitglied}
\label{beispiel-totzeitglied}

Das Totzeitglied besitzt nach Gleichung (\ref{eq:8.68}) einen Phasengang von
\begin{align}
\varphi(\Omega) &= - \Omega  k_{0} \label{eq:8.81}
\end{align}
und damit eine konstante Gruppenlaufzeit von
\begin{align}
T_{G}(\Omega) &= - \frac{d}{d\Omega}\left( - \Omega  k_{0} \right)  T_{A} = k_{0}  T_{A} \label{eq:8.82}
\end{align}
Die konstante Gruppenlaufzeit führt auch bei der Anregung eines LTI-Systems mit zwei harmonischen Signalen unterschiedlicher Frequenz zu einer identischen Verzögerung der Signale.

\end{example}

Ist die Gruppenlaufzeit eines Systems konstant, findet keine Phasenverzerrung des Systems statt. Signalanteile aller Frequenzen laufen gleich schnell durch das System. Wie die Beispiele zeigen, wird diese Bedingung nicht von allen Systemen ideal erfüllt. Es zeigt sich, dass eine weitgehend unverzerrte Übertragung eines Signals gewährleistet ist, wenn die Gruppenlaufzeit \(T_{G}\) innerhalb des Frequenzbereichs mit großer Amplitude konstant sind.

\subsection{Beispiele für FIR-Systeme mit linearer und nichtlinearer Phase}\label{sec:8.4.5}

Das in Abschnitt \ref{sec:8.4.4} diskutierte rekursive Filter besitzt keinen linearen Phasengang. Es kann gezeigt werden, dass keine rekursiven Filter (IIR-Filter) realisiert werden können, die über einen großen Frequenzbereich eine lineare Phase beziehungsweise eine konstante Gruppenlaufzeit aufweisen. An zwei Beispielen wird gezeigt, dass FIR-Systeme unter noch zu definierenden Bedingungen einen linearen Phasengang haben.

\subsubsection{FIR-System mit linearer Phase}\label{sec:fir-system-mit-linearer-phase}

Als Beispiel wird ein FIR-System mit der Übertragungsfunktion
\begin{align}
G_{1}(z) &= \frac{1}{9}  \frac{z^{4} + 2  z^{3} + 3  z^{2} + 2  z^{1} + 1}{z^{4}} = \frac{1}{9} + \frac{2}{9}  z^{- 1} + \frac{3}{9}  z^{- 2} + \frac{2}{9}  z^{- 3} + \frac{1}{9}  z^{- 4} \label{eq:8.83}
\end{align}
diskutiert. Das System hat ausschließlich Pole im Koordinatenursprung und Nullstellen auf dem Einheitskreis. Die Impulsantwort ergibt sich mit der Verschiebungsregel zu
\begin{align}
g_{1}[ k] &= \frac{1}{9}  \delta[ k] + \frac{2}{9}  \delta[ k - 1] + \frac{3}{9}  \delta[ k - 2] + \frac{2}{9}  \delta[ k - 3] + \frac{1}{9}  \delta[ k - 4] \label{eq:8.84}
\end{align}
Die Pol-Nullstellenverteilung des Systems sowie seine Impulsantwort sind in Bild \ref{fig:8.15} dargestellt.

\begin{figure}
  \centering
  \includegraphics[alt={Pol-Nullstellenverteilungen sowie die Impulsantwort des FIR-Systems \(G_{1}(z)\)}]{./media/9_8_polnullstellenverteilungen_sowie.png}
  \caption{Pol-Nullstellenverteilungen sowie die Impulsantwort des FIR-Systems \(G_{1}(z)\)}
  \label{fig:8.15}
\end{figure}

Das System weist ausschließlich Pole im Koordinatenursprung auf und ist damit stabil. Der Frequenzgang ergibt sich damit aus der Übertragungsfunktion \(G(z)\) durch die Substitution z = \(e^{jΩ}\) zu
\begin{align}
G_{1}(\Omega) &= \frac{1}{9} + \frac{2}{9}  e^{- j  \Omega} + \frac{3}{9}  e^{- j  2  \Omega} + \frac{2}{9}  e^{- j  3  \Omega} + \frac{1}{9}  e^{- j  4  \Omega} \label{eq:8.85}
\end{align}
Statt der Aufteilung in einzelne Linearfaktoren kann es bei FIR-Systemen vorteilhaft sein, den Frequenzgang \(G(Ω)\) durch Umformungen in Betrag und Phase aufzuteilen. Dazu wird die mittlere Verzögerung des Filters, im Beispiel dieses FIR-Systems \(e^{-j2Ω}\), ausgeklammert und die dadurch symmetrisch werdenden Exponentialfunktionen mit der Eulerschen Formel als Winkelfunktionen ausgedrückt. Für das System ergibt sich
\begin{align}
{G{}_{1} \left(\Omega \right)} &=  {\frac{1}{9} +\frac{2}{9}  e^{-j \Omega } +\frac{3}{9}  e^{-j 2 \Omega } +\frac{2}{9}  e^{-j 3 \Omega } +\frac{1}{9}  e^{-j 4 \Omega } } \\
{} &=  {e^{-j 2 \Omega }  \left(\frac{1}{9}  e^{j 2 \Omega } +\frac{2}{9}  e^{j \Omega } +\frac{3}{9} +\frac{2}{9}  e^{-j \Omega } +\frac{1}{9}  e^{-j 2 \Omega } \right)} \\
{} &=  {e^{-j 2 \Omega }  \left(\frac{1}{9}  e^{j 2 \Omega } +\frac{1}{9}  e^{-j 2 \Omega } +\frac{2}{9}  e^{j \Omega } +\frac{2}{9}  e^{-j \Omega } +\frac{3}{9} \right)} \\
{} &=  {e^{-j 2 \Omega }  \left(\frac{2}{9}  \cos \left(2 \Omega \right)+\frac{4}{9}  \cos \left(\Omega \right)+\frac{3}{9} \right)}
 \label{eq:8.86}
\end{align}
Der Klammerausdruck ist für alle Frequenzen im Bereich - π \textless{} Ω \textless{} π größer als null und hat demnach eine konstante Phase ϕ = 0. Da der Betrag der Exponentialfunktion mit imaginärem Argument eins ist, ergeben sich Amplituden- und Phasengang dieses FIR-Systems zu
\begin{align}
A_{1}(\Omega) &= \frac{2}{9}  \cos(2  \Omega) + \frac{4}{9}  \cos(\Omega) + \frac{3}{9} \label{eq:8.87}
\end{align}
beziehungsweise
\begin{align}
\varphi_{1}(\Omega) &= - 2  \Omega \label{eq:8.88}
\end{align}
Amplituden- und Phasengang sind in Bild \ref{fig:8.16} dargestellt.

\begin{figure}
  \centering
  \includegraphics[alt={Amplituden- und Phasengang des FIR-Systems}]{./media/9_9_amplituden_und_phasengang_des_fi.png}
  \caption{Amplituden- und Phasengang des FIR-Systems}
  \label{fig:8.16}
\end{figure}

Bei diesem FIR-System handelt es sich um ein Tiefpass-Filter mit einer linearen Phase.

\subsubsection{FIR-System mit nichtlinearer Phase}\label{sec:fir-system-mit-nichtlinearer-phase}

Nicht jedes FIR-System ist ein System mit linearer Phase. Als Beispiel für ein System mit nichtlinearer Phase wird ein System diskutiert, das eine Übertragungsfunktion von
\begin{align}
G_{2}(z) &= \frac{1}{6}  \frac{3  z^{2} + 2  z + 1}{z^{2}} = \frac{1}{6}  \left( 3 + 2  z^{- 1} + z^{- 2} \right) \label{eq:8.89}
\end{align}
besitzt. Wie jedes FIR-System hat auch dieses System ausschließlich Pole im Koordinatenursprung. Im Gegensatz zu dem ersten System liegen aber die Nullstellen nicht auf dem Einheitskreis, sondern im Inneren des Einheitskreises. Das System besitzt eine endliche Impulsantwort von
\begin{align}
g_{2}[ k] &= \frac{1}{6}  \left( 3  \delta[ k] + 2  \delta[ k - 1] + \delta[ k - 2] \right) \label{eq:8.90}
\end{align}
Die Pol-Nullstellenverteilung des Systems sowie seine Impulsantwort sind in Bild \ref{fig:8.17} dargestellt.

\begin{figure}
  \centering
  \includegraphics[alt={Pol-Nullstellenverteilungen sowie die Impulsantwort des FIR-Systems}]{./media/9_10_polnullstellenverteilungen_sowi.png}
  \caption{Pol-Nullstellenverteilungen sowie die Impulsantwort des FIR-Systems}
  \label{fig:8.17}
\end{figure}

Amplituden- und Phasengang ergeben sich nach den Rechenregeln zur Fourier-Transformation zu
\begin{align}
G_{2}(\Omega) &= \frac{1}{6}  \left( 3 + 2  e^{- j  \Omega} + e^{- j  2  \Omega} \right) \label{eq:8.91}
\end{align}
Die Koeffizienten der Impulsantwort sind unsymmetrisch. Der Frequenzgang kann deshalb mathematisch nicht wesentlich vereinfacht werden. Amplituden- und Phasengang sind in Bild \ref{fig:8.18} dargestellt.

\begin{figure}
  \centering
  \includegraphics[alt={Amplituden- und Phasengang eines unsymmetrischen FIR-Filters}]{./media/9_11_amplituden_und_phasengang_eines.png}
  \caption{Amplituden- und Phasengang eines unsymmetrischen FIR-Filters}
  \label{fig:8.18}
\end{figure}

Offensichtlich existieren FIR-Systeme mit linearer und nichtlinearer Phase, sodass sich die Frage stellt, wann ein System die Forderung nach linearer Phase erfüllt.

\subsection{Bedingungen an FIR-Systeme für einen linearen Phasengang}\label{sec:8.4.6}

Die Ausführungen in Abschnitt \ref{sec:8.4.5} zeigen, dass einige FIR-Systeme eine lineare Phase aufweisen, während andere FIR-Systeme eine nichtlineare Phase besitzen. Aus den Beispielen ergibt sich damit die Frage, welche Bedingung für einen linearen Phasengang erfüllt sein muss.

\subsubsection{Phasengang von Systemen mit Nullstellen auf dem Einheitskreis}\label{sec:phasengang-von-systemen-mit-nullstellen-auf-dem-einheitskreis}

Die Forderung nach einer linearen Phase kann in eine Symmetrieforderung für die Koeffizienten der Übertragungsfunktion von FIR-Systemen umgewandelt werden. Zu diesem Zweck wird exemplarisch eine Übertragungsfunktion mit achsensymmetrischen Koeffizienten gerader Ordnung diskutiert, bei der sich alle Nullstellen der Übertragungsfunktion auf dem Einheitskreis befinden. Es wird gezeigt, dass diese Systeme einen linearen Phasengang haben. Das Ergebnis lässt sich auf weitere Konstellationen verallgemeinern.

Im Fall einer geraden Filterordnung N und einer achsensymmetrischen Impulsantwort gilt die Beziehung
\begin{align}
g_{1}[ k] &= g_{1}[ N - k] \label{eq:8.92}
\end{align}
Für Filter mit endlicher Impulsantwort ergibt sich die Übertragungsfunktion direkt aus der Impulsantwort zu
\begin{align}
{G_{1} \left(z\right)} &=  {\sum _{k=0}^{N}g_{1} \left[k\right] z^{-k}  =\sum _{k=0}^{N/2-1}g_{1} \left[k\right] z^{-k}  +g_{1} \left[N/2\right] z^{-N/2} +\sum _{k=N/2+1}^{N}g_{1} \left[k\right] z^{-k}  } \\
{} &=  {g_{1} \left[N/2\right] z^{-N/2} +\sum _{k=0}^{N/2-1}g_{1} \left[k\right] \left(z^{-k} +z^{-\left(N-k\right)} \right) } \\
{} &=  {z^{-N/2}  \left(g_{1} \left[N/2\right]+\sum _{k=0}^{N/2-1}g_{1} \left[k\right] \left(z^{N/2-k} +z^{-\left(N/2+k\right)} \right) \right)}
 \label{eq:8.93}
\end{align}
Aufgrund der Stabilität von FIR-System ergibt sich der Frequenzgang durch die Substitution z = \(e^{jΩ}\) zu
\begin{align}
{G_{1} \left(\Omega \right)} &=  {e^{-j \frac{\Omega  N}{2} }  \left(g_{1} \left[N/2\right]+\sum _{k=0}^{N/2-1}g_{1} \left[k\right] \left(e^{j \frac{\Omega  \left(N-2 k\right)}{2} } +e^{-j \frac{\Omega  \left(N-2 k\right)}{2} } \right) \right)} \\
{} &=  {e^{-j \frac{\Omega  N}{2} }  \left(g_{1} \left[N/2\right]+2 \sum _{k=0}^{N/2-1}g_{1} \left[k\right] \cos \left(\frac{\Omega  \left(N-2 k\right)}{2} \right) \right)}
 \label{eq:8.94}
\end{align}
Der zweite Faktor in Gleichung (\ref{eq:8.94}) ist reell und kann positive und negative Werte annehmen. Da er stetig in Ω ist, weist er im Punkt seines Vorzeichenwechsels einen Betrag von null auf. Damit besitzt der Frequenzgang \(G_{1}(Ω)\) eine abschnittsweise lineare Phase, die bei Vorzeichenwechseln des Frequenzgangs um π springt. An diesen Stellen ist aber der Betrag der Übertragungsfunktion null. Das Filter weist damit im gesamten Durchlassbereich eine lineare Phase auf. Als Beispiel wird ein Filter mit der Impulsantwort
\begin{align}
g_{1}[ k] &= \frac{1}{3.7}  \left( \delta[ k] + 2.7  \delta[ k - 1] + 3.7  \delta[ k - 2] + 2.7  \delta[ k - 3] + \delta[ k - 4] \right) \label{eq:8.95}
\end{align}
betrachtet. Bild \ref{fig:8.19} stellt Pol-Nullstellen-Diagramm, Impulsantwort, Amplituden- und Phasengang des Filters dar.

\begin{figure}
  \centering
  \includegraphics[alt={Pol-Nullstellen-Diagramm, Impulsantwort, Amplituden- und Phasengang\\
eines FIR-Filters mit gerader Ordnung und achsensymmetrischer Impulsantwort}]{./media/image203.png}
  \caption{Pol-Nullstellen-Diagramm, Impulsantwort, Amplituden- und Phasengang\\
eines FIR-Filters mit gerader Ordnung und achsensymmetrischer Impulsantwort}
  \label{fig:8.19}
\end{figure}

Der Amplitudengang wird für Ω = ± 2/3π und Ω = ± 5/6π zu null, und es finden Vorzeichenwechsel statt. Deshalb springt der Phasengang an diesen Stellen um π. Der Phasengang ist abschnittsweise linear.

Die an dem Beispiel gewonnenen Ergebnisse lassen sich auf weitere FIR-Systeme verallgemeinern. FIR-Systeme, bei denen die Nullstellen ausschließlich auf dem Einheitskreis liegen, weisen einen linearen Phasengang auf.

\subsubsection{Weitere Symmetriebedingungen für einen linearen Phasengang}\label{sec:weitere-symmetriebedingungen-fuxfcr-einen-linearen-phasengang}

In Abschnitt \ref{sec:8.3.2} wird der Frequenzgang zeitdiskreter Systeme mit gebrochen rationaler Übertragungsfunktion diskutiert. Er setzt sich aus den Frequenzgängen einzelner Linearfaktoren zusammen. Liegen die Nullstellen dieser Linearfaktoren auf dem Einheitskreis, sind die zugehörigen Phasengänge linear zur Frequenz Ω. Das Beispiel in Abschnitt \ref{sec:8.4.6} bestätigt diesen Zusammenhang.

Zusätzlich zu den Nullstellen auf dem Einheitskreis können Nullstellen-Konstellationen wie in Gleichung (\ref{eq:8.96}) auftreten.
\begin{align}
{G\left(z\right)} &=  {\frac{\left(z-r_{0}  e^{j \Omega _{0} } \right) \left(z-r_{0}  e^{-j \Omega _{0} } \right) \left(z-\frac{1}{r_{0} }  e^{j \Omega _{0} } \right) \left(z-\frac{1}{r_{0} }  e^{-j \Omega _{0} } \right)}{z^{4} } } \\
{} &=  {\left(1-\frac{r_{0}  e^{j \Omega _{0} } }{z} \right) \left(1-\frac{r_{0}  e^{-j \Omega _{0} } }{z} \right) \left(1-\frac{e^{j \Omega _{0} } }{r_{0}  z} \right) \left(1-\frac{e^{-j \Omega _{0} } }{r_{0}  z} \right)}
 \label{eq:8.96}
\end{align}
Die Nullstellen sind konjugiert komplex und haben einen Abstand \(r_{0}\) und 1/\(r_{0}\) vom Ursprung. Da die Nullstellen nicht auf dem Einheitskreis liegen, finden keine Phasensprünge statt. Jede Nullstelle für sich hat damit keinen linearen Phasengang. Zur Berechnung des Phasengangs der gesamten Konstellation wird zunächst der Frequenzgang des Systems bestimmt.
\begin{align}
{G\left(\Omega \right)} &=  {\left(1-\frac{r_{0}  e^{j \Omega _{0} } }{e^{j \Omega } } \right) \left(1-\frac{r_{0}  e^{-j \Omega _{0} } }{e^{j \Omega } } \right) \left(1-\frac{e^{j \Omega _{0} } }{r_{0}  e^{j \Omega } } \right) \left(1-\frac{e^{-j \Omega _{0} } }{r_{0}  e^{j \Omega } } \right)} \\
{} &=  {\left(1-r_{0}  e^{j \left(\Omega _{0} -\Omega \right)} \right) \left(1-r_{0}  e^{-j \left(\Omega _{0} +\Omega \right)} \right) \left(1-\frac{e^{j \left(\Omega _{0} -\Omega \right)} }{r_{0} } \right) \left(1-\frac{e^{-j \left(\Omega _{0} +\Omega \right)} }{r_{0} } \right)}
 \label{eq:8.97}
\end{align}
Ausmultiplizieren ergibt
\begin{align}
{G\left(\Omega \right)} &=  {\left(1-r_{0}  e^{j \left(\Omega _{0} -\Omega \right)} \right) \left(1-r_{0}  e^{-j \left(\Omega _{0} +\Omega \right)} \right) \left(1-\frac{e^{j \left(\Omega _{0} -\Omega \right)} }{r_{0} } \right) \left(1-\frac{e^{-j \left(\Omega _{0} +\Omega \right)} }{r_{0} } \right)} \\
{} &=  {\left(1-2 r_{0}  e^{-j \Omega }  \cos \left(\Omega _{0} \right)+r_{0}^{2}  e^{-j 2 \Omega } \right) \left(1-2 \frac{1}{r_{0} }  e^{-j \Omega }  \cos \left(\Omega _{0} \right)+\frac{1}{r_{0}^{2} }  e^{-j 2 \Omega } \right)} \\
{} &=  {1-2 r_{0}  e^{-j \Omega }  \cos \left(\Omega _{0} \right)+r_{0}^{2}  e^{-j 2 \Omega } } \\
{} & {} & {-2 \frac{1}{r_{0} }  e^{-j \Omega } \cos \left(\Omega _{0} \right)+4 e^{-j 2 \Omega } \cos {}^{2} \left(\Omega _{0} \right)-2 r_{0}  e{}^{-j 3 \Omega }  \cos \left(\Omega _{0} \right)} \\
{} & {} & {+\frac{1}{r_{0}^{2} }  e^{-j 2 \Omega } -2 \frac{1}{r_{0} }  e^{-j 3 \Omega } \cos \left(\Omega _{0} \right)+e^{-j 4 \Omega } }
 \label{eq:8.98}
\end{align}
Die Summanden werden so sortiert, dass Paare von Exponentialfunktionen eine Kosinusfunktion bilden.
\begin{align}
{G\left(\Omega \right)} &=  {1+e^{-j 4 \Omega } -2 r_{0}  \cos \left(\Omega _{0} \right) \left(e^{-j \Omega } +e^{-j 3 \Omega } \right)} \\
{} & {} & {+r_{0}^{2}  e{}^{-j 2 \Omega } +4 e{}^{-j 2 \Omega } \cos {}^{2} \left(\Omega _{0} \right)+\frac{1}{r_{0}^{2} }  e^{-j 2 \Omega } -2 \frac{1}{r_{0} } \cos \left(\Omega _{0} \right) \left(e^{-j \Omega } +e^{-j 3 \Omega } \right)} \\
{} &=  {e^{-j 2 \Omega }  \left(\left(e^{j 2 \Omega } +e^{-j 2 \Omega } \right)-2 r_{0}  \cos \left(\Omega _{0} \right) \left(e^{j \Omega } +e^{-j \Omega } \right)\right)} \\
{} & {} & {+e^{-j 2 \Omega }  \left(r_{0}^{2} +4 \cos ^{2} \left(\Omega _{0} \right)+\frac{1}{r_{0}^{2} } -2 \frac{1}{r_{0} }  \cos \left(\Omega _{0} \right) \left(e^{j \Omega } +e^{-j \Omega } \right)\right)} \\
{} &=  {e^{-j 2 \Omega }  \left(2 \cos \left(2 \Omega \right)-4 r_{0}  \cos \left(\Omega _{0} \right) \cos \left(\Omega \right)\right)} \\
{} & {} & {+e^{-j 2 \Omega }  \left(r_{0}^{2} +4 \cos ^{2} \left(\Omega _{0} \right)+\frac{1}{r_{0}^{2} } -4 \frac{1}{r_{0} }  \cos \left(\Omega _{0} \right) \cos \left(\Omega \right)\right)}
 \label{eq:8.99}
\end{align}
Der Faktor \(e^{-j2Ω}\) repräsentiert die lineare Phasenverschiebung. Der zweite Faktor ist reell. Um zu zeigen, dass in dem Intervall - π ≤ Ω ≤ π keine Phasensprünge auftreten, werden die Nullstellen berechnet. Unter Berücksichtigung von
\begin{align}
\cos(2  \Omega) &= 2  \cos^{2}(\Omega) - 1 \label{eq:8.100}
\end{align}
ergibt sich als Bedingung für die Nullstellen
\begin{align}
{0} &=  {\cos ^{2} \left(\Omega \right)-\left(r_{0} +\frac{1}{r_{0} } \right) \cos \left(\Omega _{0} \right) \cos \left(\Omega \right)+\frac{1}{4}  \left(r_{0}^{2} -2+\frac{1}{r_{0}^{2} } \right)+\cos ^{2} \left(\Omega _{0} \right)} \\
{} &=  {\cos ^{2} \left(\Omega \right)-\left(r_{0} +\frac{1}{r_{0} } \right) \cos \left(\Omega _{0} \right) \cos \left(\Omega \right)+\frac{1}{4}  \left(r_{0} -\frac{1}{r_{0} } \right)^{2} +\cos ^{2} \left(\Omega _{0} \right)}
 \label{eq:8.101}
\end{align}
Es handelt sich um eine quadratische Gleichung in cos(Ω). Die Lösung der quadratischen Gleichungen ergibt
\begin{align}
\cos(\Omega) &= \frac{1}{2}  \left( r_{0} + \frac{1}{r_{0}} \right)  \cos\left( \Omega_{0} \right) \pm \sqrt{\frac{1}{4}  \left( r_{0} + \frac{1}{r_{0}} \right)^{2}  \cos^{2}\left( \Omega_{0} \right) - \frac{1}{4}  \left( r_{0} - \frac{1}{r_{0}} \right)^{2} + \cos^{2}\left( \Omega_{0} \right)} \label{eq:8.102}
\end{align}
Das Argument der Wurzel ist negativ.
\begin{align}
{\frac{1}{4}  \left( r_{0} + \frac{1}{r_{0}} \right)^{2}  \cos^{2}\left( \Omega_{0} \right) - \frac{1}{4}  \left( r_{0} - \frac{1}{r_{0}} \right)^{2} - \cos^{2}\left( \Omega_{0} \right)
}{= \frac{1}{4}  \left( r_{0}^{2} + 2 + \frac{1}{r_{0}^{2}} \right)  \cos^{2}\left( \Omega_{0} \right) - \frac{1}{4}  \left( r_{0}^{2} - 2 + \frac{1}{r_{0}^{2}} \right)^{2} - \cos^{2}\left( \Omega_{0} \right)
}{= \frac{1}{4}  \left( r_{0}^{2} + \frac{1}{r_{0}^{2}} \right)  \cos^{2}\left( \Omega_{0} \right) + \frac{1}{2}  \cos^{2}\left( \Omega_{0} \right) - \frac{1}{4}  \left( r_{0}^{2} + \frac{1}{r_{0}^{2}} \right)^{2} + \frac{1}{2} - \cos^{2}\left( \Omega_{0} \right)
}{= \frac{1}{4}  \left( r_{0}^{2} + \frac{1}{r_{0}^{2}} \right)  \left( \cos^{2}\left( \Omega_{0} \right) - 1 \right) + \frac{1}{2}  \left( 1 - \cos^{2}\left( \Omega_{0} \right) \right) = - \frac{1}{4}  \left( r_{0} - \frac{1}{r_{0}} \right)^{2}  \sin^{2}\left( \Omega_{0} \right) < 0} \label{eq:8.103}
\end{align}
Es existieren also keine reellen Nullstellen, sodass der Ausdruck nach der Exponentialfunktion keine Vorzeichenwechsel aufweist. Nullstellen-Konstellationen wie in Gleichung (\ref{eq:8.96}) besitzen damit eine lineare Phase. In Bild \ref{fig:8.20} sind Impulsantwort, Pol-Nullstellen-Diagramm, Amplituden- und Phasengang für ein Beispiel mit \(r_{0} = 0.75\) und \(Ω_{0} = 2\)/3π dargestellt.

\begin{figure}
  \centering
  \includegraphics[alt={Pol-Nullstellen-Diagramm, Impulsantwort, Amplituden- und Phasengang\\
eines FIR-Filters mit symmetrischer Lage von Nullstellen um den Einheitskreis}]{./media/image209.png}
  \caption{Pol-Nullstellen-Diagramm, Impulsantwort, Amplituden- und Phasengang\\
eines FIR-Filters mit symmetrischer Lage von Nullstellen um den Einheitskreis}
  \label{fig:8.20}
\end{figure}

Die numerische Berechnung bestätigt, dass die mit Gleichung (\ref{eq:8.96}) beschriebene Nullstellenkonstellation ebenfalls zu einem System mit linearem Phasengang führt.

\subsubsection{\texorpdfstring{Zusammenfassung von Bedingungen an FIR-Systeme mit linearem Phasengang }{Zusammenfassung von Bedingungen an FIR-Systeme mit linearem Phasengang }}\label{zusammenfassung-von-bedingungen-an-fir-systeme-mit-linearem-phasengang}

Tabelle \ref{tab:8.1} fasst die Bedingungen für die Pol-Nullstellen für FIR-Filter mit linearer Phase zusammen, die in diesem Kapitel behandelt werden.

\begin{table}
  \centering
  \caption{Bedingungen an die Pol-Nullstellen-Lagen für FIR-Filter mit linearer Phase}
  \label{tab:8.1}
  \begin{tabular}{ll}
    \toprule
    Alle Pole an der Stelle & \(\alpha = 0\) \\
    \midrule
    Reelle Nullstellen auf dem Einheitskreis & \(\beta = 0und\beta = \pi\) \\
    Konjugiert komplexe Nullstellen 
auf dem Einheitskreis
 & \(\beta_{1,2} = e^{\pm j  \Omega_{0}}\) \\
    Nullstellenkonstellation 
symmetrisch zum Einheitskreis
 & \(\beta_{1,2} = r_{0}  e^{\pm j  \Omega_{0}}und\beta_{3,4} = \frac{1}{r_{0}}  e^{\pm j  \Omega_{0}}\) \\
    \bottomrule
  \end{tabular}
\end{table}

\subsection{Vergleich von Tiefpassfiltern mit linearer und nichtlinearer Phase}\label{sec:8.4.7}

Die Diskussion der Beispiele im letzten Abschnitt zeigt, dass Systeme mit linearem Phasengang Eingangssignale zwar verzögern, jedoch nicht in ihrer Form verzerren. Im Folgenden wird analysiert, welcher Unterschied zwischen zwei Tiefpass-Filtern besteht, die eine lineare beziehungsweise eine nichtlineare Phase besitzen. Das Filter mit linearer Phase wird als FIR-Filter realisiert, das Filter mit nichtlinearer Phase als IIR-Filter. Beide Filter werden so entworfen, dass sie bei einer Frequenz von \(Ω = π\)/5 eine Dämpfung von maximal 1~dB aufweisen und bei einer Frequenz von \(Ω = 1\).1π/5 eine Dämpfung von mindestens 40~dB besitzen. Die Forderung führt zu Filtern hoher Ordnung, an denen die Unterschiede zwischen linearem und nichtlinearem Phasengang deutlicher zu erkennen sind. Amplituden- und Phasengang der beiden Filter sind in Bild \ref{fig:8.21} dargestellt.

\begin{figure}
  \centering
  \includegraphics[alt={Frequenzgangskennlinien des FIR und IIR-Filters}]{./media/9_12_frequenzgangskennlinien_des_fir.png}
  \caption{Frequenzgangskennlinien des FIR und IIR-Filters}
  \label{fig:8.21}
\end{figure}

Die Amplitudengänge entsprechen den Vorgaben aus dem Filterdesign. Die Phasengänge weichen stark voneinander ab. Das FIR-Filter hat einen linearen Phasenverlauf, dem an den Nullstellen des Amplitudengangs zusätzlich Phasensprünge überlagert sind. Im Gegensatz dazu hat das IIR-Filter einen nichtlinearen Phasenverlauf, dem ebenfalls Phasensprünge überlagert sind. Auch die Absolutwerte der Phase sind stark unterschiedlich.

Zunächst wird analysiert, wie zwei Wellenpakete mit unterschiedlicher Frequenz \(Ω_{1} = π\)/50 und \(Ω_{2}\)~=~π/15 das Filter durchlaufen. Bild \ref{fig:8.22} zeigt die Systemreaktion der Tiefpass-Filter auf die Anregung.

\begin{figure}
  \centering
  \includegraphics[alt={Reaktion von FIR-Filter mit linearer Phase und IIR-Filter mit nichtlinearer Phase\\
auf die Anregung mit zwei Wellenpakten unterschiedlicher Frequenz}]{./media/image211.png}
  \caption{Reaktion von FIR-Filter mit linearer Phase und IIR-Filter mit nichtlinearer Phase\\
auf die Anregung mit zwei Wellenpakten unterschiedlicher Frequenz}
  \label{fig:8.22}
\end{figure}

Die beiden Wellenpakete werden mit einem Zeitabstand von 500~Abtastzeiten gesendet. Die Verzögerung der Wellenpakete entspricht der Gruppenlaufzeit. Sie ist wegen der linearen Phase unabhängig von der Frequenz und ergibt sich aus der Filterordnung. In diesem Beispiel beträgt die Filterordnung\\
\(K = 201\) und die Verzögerung ist
\begin{align}
\frac{T_{G}}{T_{A}} &= k_{G} = \frac{K}{2} = 100.5 \label{eq:8.104}
\end{align}
Das System mit linearer Phase weist eine große Gruppenlaufzeit auf, die gleichbedeutend mit einer großen Phase ist. Bei dem System mit nichtlinearer Phase ist insbesondere das zweite Wellenpaket verzerrt. Bild \ref{fig:8.23} zeigt die Systemreaktion des FIR- und IIR-Filters auf die Anregung mit Rechteckfolgen.

\begin{figure}
  \centering
  \includegraphics[alt={Reaktion von FIR-Filter mit linearer Phase und IIR-Filter mit nichtlinearer Phase\\
auf die Anregung mit Rechteckfolgen}]{./media/image212.png}
  \caption{Reaktion von FIR-Filter mit linearer Phase und IIR-Filter mit nichtlinearer Phase\\
auf die Anregung mit Rechteckfolgen}
  \label{fig:8.23}
\end{figure}

Deutlich zu erkennen sind die größere Gruppenlaufzeit und die geringeren Verzerrungen des FIR-Filters mit linearer Phase. Das IIR-Filter mit nichtlinearer Phase besitzt eine kürzere Gruppenlaufzeit, aber das Ausgangssignal ist unsymmetrisch und damit stärker verzerrt. Da sich ein Signal im allgemeinen Fall aus vielen Spektralanteilen zusammensetzt, ist für die verzerrungsarme Übertragung ein Filter mit linearer Phase im Durchlassbereich notwendig.

Aus diesen Überlegungen ergeben sich zwei Forderungen an zeitdiskrete Filter. Sie müssen eine minimale Verzögerung besitzen und im Durchlassbereich einen zumindest näherungsweise linearen Phasengang aufweisen.

\section{\texorpdfstring{Allpässe und minimalphasige Systeme }{Allpässe und minimalphasige Systeme }}\label{allpuxe4sse-und-minimalphasige-systeme}

In den Abschnitten \ref{sec:8.3} und \ref{sec:8.4} wird die Bedeutung der Phase eines Systems diskutiert. Es zeigt sich, dass eine große Phase einer großen Zeitverzögerung entspricht und umgekehrt. In der Regelungstechnik und Signalverarbeitung ist es gewünscht, das Ausgangssignal mit einer möglichst kleinen Verzögerung auszugeben. Das entsprechende System muss dazu eine minimale Phase aufweisen.

Systeme sind durch die Angabe ihres Amplitudengangs nicht eindeutig definiert. Aus der Forderung nach Stabilität folgt, dass die Pole der Übertragungsfunktion im Einheitskreis liegen müssen. Bei der Diskussion der Invertierbarkeit des Systems in Kapitel \ref{sec:6.3.9} wird gefordert, dass auch die Nullstellen der Übertragungsfunktion im Einheitskreis liegen müssen. Diese Forderung führt zu sogenannten minimalphasigen Systemen. Sie zeichnen sich durch eine Systemreaktion mit möglichst geringer Verzögerung aus.

Manche Systeme sind nicht minimalphasig, können aber in einen minimalphasigen Anteil und einen Allpass zerlegt werden.

\subsection{Minimalphasige Systeme}\label{sec:minimalphasige-systeme}

Zur Veranschaulichung des Begriffes der Minimalphasigkeit werden zwei Systeme miteinander verglichen, die denselben Amplitudengang aber einen unterschiedlichen Phasengang besitzen. Sie haben die Übertragungsfunktionen
\begin{align}
G_{1}(z) &= \frac{(z - 2  j)  (z + 2  j)}{z^{2}} = \frac{z^{2} + 4}{z^{2}} \label{eq:8.105}
\end{align}
und
\begin{align}
G_{2}(z) &= 8  \frac{(z - 0.5  j)  (z + 0.5  j)}{z^{2}} = 4  \frac{z^{2} + 0.25}{z^{2}} \label{eq:8.106}
\end{align}
Ihre Pol-Nullstellen-Diagramme sind in Bild \ref{fig:8.24}, Amplituden- und Phasengänge sind in Bild \ref{fig:8.25} dargestellt.

\begin{figure}
  \centering
  \includegraphics[alt={Pol-Nullstellen-Diagramme, Amplituden- und Phasengang von zwei Übertragungssystemen}]{./media/9_13_polnullstellendiagramme_amplitu.png}
  \caption{Pol-Nullstellen-Diagramme, Amplituden- und Phasengang von zwei Übertragungssystemen}
  \label{fig:8.24}
\end{figure}

\begin{figure}
  \centering
  \includegraphics[alt={Pol-Nullstellen-Diagramme, Amplituden- und Phasengang von zwei Übertragungssystemen}]{./media/9_14_polnullstellendiagramme_amplitu.png}
  \caption{Pol-Nullstellen-Diagramme, Amplituden- und Phasengang von zwei Übertragungssystemen}
  \label{fig:8.25}
\end{figure}

Die Systeme haben denselben Amplitudengang \(A(Ω)\). Ihre Phasengänge ϕ(Ω) unterscheiden sich jedoch erheblich. Die unterschiedlichen Phasengänge führen erwartungsgemäß zu unterschiedlichen Signalverzögerungen, die zum Beispiel an den beiden Sprungantworten in Bild \ref{fig:8.26} abgelesen werden können.

\begin{figure}
  \centering
  \includegraphics[alt={Sprungantworten der diskutierten Übertragungssysteme}]{./media/9_15_sprungantworten_der_diskutierte.png}
  \caption{Sprungantworten der diskutierten Übertragungssysteme}
  \label{fig:8.26}
\end{figure}

Das in diesem Beispiel dargestellte Verhalten kann verallgemeinert werden. Systeme, die ausschließlich Nullstellen innerhalb des Einheitskreises besitzen, besitzen eine minimale Phase. Sie werden als minimalphasige Systeme bezeichnet.

Die Forderung nach Nullstellen innerhalb des Einheitskreises entspricht der Forderung nach Invertierbarkeit, sodass alle invertierbaren Systeme eine minimale Phase aufweisen und alle minimalphasigen Systeme invertierbar sind.

\subsection{Allpässe}\label{sec:allpuxe4sse}

Übertragungsglieder, die Spektralanteile aller Frequenzen mit dem gleichen Betrag übertragen, werden als Allpass bezeichnet. Sie weisen einen Amplitudengang von
\begin{align}
A(\Omega) &= \left| G(\Omega) \right| = 1 \label{eq:8.107}
\end{align}
auf. Ein Beispiel für einen Allpass ist ein System mit einer Verzögerung um \(k_{0}\).
\begin{align}
G(\Omega) &= e^{- j  k_{0}  \Omega} \label{eq:8.108}
\end{align}
Der Frequenzgang des Systems mit zeitlicher Verzögerung hat einen konstanten Betrag von \(A(Ω)\) = 1 und eine Phase von ϕ(Ω) = - k\textsubscript{0}Ω. Allpässe wirken sich damit nur auf die Phase des Eingangssignals aus. Allpässe mit gebrochen rationalen Übertragungsfunktionen besitzen Nullstellen außerhalb des Einheitskreises. Sie haben die Struktur
\begin{align}
G(z) &= \frac{z^{- 1} - \alpha^{*}}{1 - \alpha  z^{- 1}} = \frac{1 - \alpha^{*}  z}{z - \alpha} \label{eq:8.109}
\end{align}
Beim Allpass gehört zu jedem Pol innerhalb des Einheitskreises eine Nullstelle außerhalb des Einheitskreises mit gleicher Phase \(Ω_{0}\) und reziprokem Betrag 1/\(r_{0}\). Sämtliche Nullstellen eines stabilen Allpass liegen außerhalb des Einheitskreises. Bild \ref{fig:8.27} stellt die Lage des Pols und der Nullstelle an einem Beispiel dar.

\begin{figure}
  \centering
  \includegraphics[alt={Lage von Pol und Nullstelle eines Allpasses erster Ordnung}]{./media/9_16_lage_von_pol_und_nullstelle_ein.png}
  \caption{Lage von Pol und Nullstelle eines Allpasses erster Ordnung}
  \label{fig:8.27}
\end{figure}

Der Frequenzgang ergibt sich mit
\begin{align}
\alpha &= r_{0}  e^{j  \Omega_{0}} \label{eq:8.110}
\end{align}
zu
\begin{align}
G(\Omega) &= \frac{e^{- j  \Omega} - r_{0}  e^{- j  \Omega_{0}}}{1 - r_{0}  e^{j  \Omega_{0}}  e^{- j  \Omega}} = e^{- j  \Omega}  \frac{1 - r_{0}  e^{j  \left( \Omega - \Omega_{0} \right)}}{1 - r_{0}  e^{- j  \left( \Omega - \Omega_{0} \right)}} \label{eq:8.111}
\end{align}
Zähler und Nenner des Frequenzgangs sind konjugiert komplex zueinander, sie haben demnach denselben Betrag. Die Exponentialfunktion mit imaginärem Argument hat ebenfalls den Betrag 1. Der Amplitudengang \(A(Ω)\) hat damit den konstanten Betrag von \(A = 1\). Der Phasengang ergibt sich zu
\begin{align}
{\varphi \left(\Omega \right)} &=  {-\Omega +\arctan \left(\frac{r_{0}  \sin \left(\Omega _{0} -\Omega \right)}{1-r_{0}  \cos \left(\Omega _{0} -\Omega \right)} \right)-\arctan \left(\frac{-r_{0}  \sin \left(\Omega _{0} -\Omega \right)}{1-r_{0}  \cos \left(\Omega _{0} -\Omega \right)} \right)} \\
{} &=  {-\Omega +2 \arctan \left(\frac{r_{0}  \sin \left(\Omega _{0} -\Omega \right)}{1-r_{0}  \cos \left(\Omega _{0} -\Omega \right)} \right)}
 \label{eq:8.112}
\end{align}
Amplitudengang und Phasengang des Allpasses mit Pol- und Nullstelle aus Bild \ref{fig:8.27} sind in Bild \ref{fig:8.28} dargestellt.

\begin{figure}
  \centering
  \includegraphics[alt={Amplituden- und Phasengang eines Allpasses erster Ordnung}]{./media/9_17_amplituden_und_phasengang_eines.png}
  \caption{Amplituden- und Phasengang eines Allpasses erster Ordnung}
  \label{fig:8.28}
\end{figure}

Bild \ref{fig:8.28} zeigt, dass Allpässe nur die Phasenlage verändern und die Amplituden unverändert lassen.

\subsection{Entwurf von minimalphasigen Systemen}\label{sec:entwurf-von-minimalphasigen-systemen}

Zur Transformation eines nichtminimalphasigen Systems in ein minimalphasiges System wird das System mit einem mit einem Allpass erweitert, der die Nullstellen außerhalb des Einheitskreises kompensiert und durch Nullstellen im Einheitskreis ersetzt. Weist ein System \(G(z)\) außerhalb des Einheitskreises die Nullstellen \(β_{1}\) \ldots{} \(β_{M}\) auf, wird das System mit einem Allpass
\begin{align}
G_{A}(z) &= \prod_{m = 1}^{M}\frac{z^{- 1} - \beta_{m}^{*}}{1 - \beta_{m}  z^{- 1}} = \prod_{i = 1}^{M}\frac{1 - \beta_{m}^{*}  z}{z - \beta_{m}} \label{eq:8.113}
\end{align}
gefiltert. Da die Nullstellen außerhalb des Einheitskreises durch die Pole der Allpässe kompensiert werden, besitzt das System \(G_{M}(z)\)
\begin{align}
G_{M}(z) &= G(z)  G_{A}(z) \label{eq:8.114}
\end{align}
lediglich Nullstellen, die innerhalb des Einheitskreises liegen. Das System \(G_{M}(z)\) ist damit minimalphasig.

Zur Kompensation der Nullstellen außerhalb des Einheitskreises wird ein System eingesetzt, dass Pole außerhalb des Einheitskreises besitzt. Das resultierende System ist nur dann stabil, wenn sich Pole und Nullstellen außerhalb des Einheitskreises perfekt kompensieren. Ansonsten entsteht ein System mit Polen außerhalb des Einheitskreises, das instabil ist. Besondere Vorsicht ist bei Systemen geboten, die mit Fixed-Point-Format arbeiten. Diese ideale Kompensation wird zum Beispiel bei Kammfiltern realisiert {[}Oppe04{]}.

\begin{example}{Minimalphasiges FIR-Filter}
\label{beispiel-minimalphasiges-fir-filter}

Als Beispiel wird ein FIR-Filter mit der Übertragungsfunktion
\begin{align}
G(z) &= \frac{(z - 2)  (2  z + 1)  \left( z^{2} - 2  z + 2 \right)}{z^{4}} \label{eq:8.115}
\end{align}
diskutiert. Es hat Nullstellen bei \(β_{1} = 2\) und \(β_{2}\) = -1/2. Die beiden anderen Nullstellen berechnen sich aus
{\def\LTcaptype{none} % do not increment counter
\begin{center}
  \captionof{table}{Tabellarische Übersicht}
  \begin{tabular}{l}
    \toprule
    \noalign{}

\[\beta_{3,4} = - \frac{- 2}{2} \pm \sqrt{1 - 2} = 1 \pm j = \sqrt{2}  e^{\pm j  \frac{\pi}{4}}\] \\
    \bottomrule
  \end{tabular}
\end{center}
}Damit das System minimalphasig wird, müssen alle Nullstellen innerhalb des Einheitskreises liegen. Dazu wird die Übertragungsfunktion mit der Übertragungsfunktion des Allpass
\begin{align}
G_{A}(z) &= \frac{1 - 2  z}{z - 2}  \frac{1 - \sqrt{2}  e^{- j  \frac{\pi}{4}}  z}{z - \sqrt{2}  e^{j  \frac{\pi}{4}}}  \frac{1 - \sqrt{2}  e^{j  \frac{\pi}{4}}  z}{z - \sqrt{2}  e^{- j  \frac{\pi}{4}}} = \frac{(1 - 2  z)  \left( 1 - 2  z - 2  z^{2} \right)}{(z - 2)  \left( z^{2} - 2  z + 2 \right)} \label{eq:8.116}
\end{align}
multipliziert.
\begin{align}
{G_{M} \left(z\right)} &=  {\frac{\left(z-2\right) \left(2 z+1\right) \left(z^{2} -2 z+2\right)}{z^{4} }  \frac{\left(1-2 z\right) \left(1-2 z-2 z^{2} \right)}{\left(z-2\right) \left(z^{2} -2 z+2\right)} } \\
{} &=  {\frac{\left(2 z+1\right) \left(1-2 z\right) \left(1-2 z-2 z^{2} \right)}{z^{4} } }
 \label{eq:8.117}
\end{align}
Das minimalphasige System \(G_{M}(z)\) besitzt die Nullstellen bei \(β_{1,2}\) = ± 1/2 und an den Stellen\\
\(β_{3,4} = 1\)/2 ± j1/2. Alle Nullstellen liegen innerhalb des Einheitskreises.

\end{example}
